\documentclass[%
 reprint,
 amsmath,amssymb,
 aps,floatfix
]{revtex4-2}
\usepackage{subfigure}
\usepackage{physics}
\usepackage{braket}
\usepackage{lipsum} 
 \usepackage[colorlinks=true]{hyperref}
\usepackage{orcidlink}
\usepackage{graphicx}% Include figure files
\usepackage{dcolumn}% Align table columns on decimal point
\usepackage{bm}\allowdisplaybreaks

\usepackage[normalem]{ulem}
\usepackage{dsfont}
\usepackage{amsmath}
\usepackage{amssymb}
\usepackage{mathrsfs}

\usepackage{bibentry}



\begin{document}


\title{Bosonization of Non-Fermi Liquids}

\author{SangEun Han\,\orcidlink{0000-0003-3141-1964}}
\affiliation{Department of Physics, University of Toronto, Toronto, Ontario M5S 1A7, Canada}
\author{F\'{e}lix Desrochers\,\orcidlink{0000-0003-1211-901X}}
\affiliation{Department of Physics, University of Toronto, Toronto, Ontario M5S 1A7, Canada}
\author{Yong Baek Kim}
\affiliation{Department of Physics, University of Toronto, Toronto, Ontario M5S 1A7, Canada}

\date{\today}

\begin{abstract}
Understanding non-Fermi liquids in dimensions higher than one remains one of the most formidable challenges in modern condensed matter physics. These systems, characterized by an abundance of gapless degrees of freedom and the absence of well-defined quasiparticles, defy conventional analytical frameworks. Inspired by recent work on the bosonization of Fermi surfaces [Delacretaz, Du, Mehta, and Son, \href{https://doi.org/10.1103/PhysRevResearch.4.033131}{Physical Review Research, 4, 033131 (2022)}], we present a procedure for bosonizing non-Fermi liquids, which provides a holistic approach capable of addressing their intricate physics. Our method involves parameterizing the generalized fermionic distribution function through a bosonic field that describes frequency-dependent local variations of the chemical potential in momentum space. We propose an effective action that produces the collisionless quantum Boltzmann equation as its equation of motion and can straightforwardly be used for any dimension and Fermi surface of interest. Even at the quadratic order, this action reproduces highly non-trivial results obtainable only through involved analysis with alternative means. By offering a comprehensive description of the physics of non-Fermi liquids, our work stands as an important building block in advancing the comprehension of strange metals and associated phenomena.
\end{abstract}


\maketitle

%%%%%%%%%%%%%%%%%%%%%%%%%%%%%%%%%%%%%%%%%%%%%%%%%%%%%%%%%%%%%%%%%%%%%%%%%%%
\section{Introduction}

The description of metals is arguably the oldest problem in all of condensed matter physics~\cite{sommerfeld1928elektronentheorie,bloch1929quantenmechanik}. Despite this historical status, a satisfactory theoretical description of the low-energy physics of metals remains an outstanding problem made exceedingly hard by the continuum of gapless degrees of freedom located along a nontrivial submanifold of momentum space (i.e., the Fermi surface). When metals host well-defined quasiparticles, Landau's phenomenological Fermi liquid (FL) theory can be used as a bare-bones description to compute certain correlation functions at leading order. Nonetheless, in many cases of physical interest, interactions are so important that the system does not have well-defined quasiparticles while maintaining a sharp Fermi surface. These so-called non-Fermi liquids (NFL) are ubiquitous in the phase diagrams of strongly correlated systems. They can arise when metals are subjected to singular interactions mediated by critical bosonic modes~\cite{altshuler1994low, kim1994gauge, Kim1995, polchinski1994low, lee1992gauge, oganesyan2001quantum, lee2008stability, fitzpatrick2014non, abanov2000spin, moon2013non, han2019emergent, nayak1994non, nayak1994renormalization, lee2009low, metlitski2010quantum1, metlitski2010quantum2, mross2010controlled, dalidovich2013perturbative, lee2018recent, han2022non}, at quantum critical points associated with a change in volume
of the Fermi surface~\cite{coleman2001fermi, han2022non2, han2022microscopic, senthil2004weak, senthil2008critical, senthil2008theory}, or when fermions interact through random couplings~\cite{sachdev1993gapless, sachdev2010holographic, maldacena2016remarks, chowdhury2022sachdev}. As they are in a strong coupling regime, constructing a controlled general theoretical framework to describe NFL has thus far remained an enticing but challenging objective despite intense interest.

Considering their fundamental role in strongly correlated electron systems, there is a crucial need for a more systematic understanding of NFL beyond perturbative treatments in more than one dimension. Substantial recent progress has been made in that regard. In particular, some universal constraints on the low-energy physics of metals --- that could apply for both FL and NFL --- have been derived using non-perturbative arguments coming from emergent symmetry groups and their corresponding anomalies~\cite{shi2022gifts, else2021non, else2021strange, lu2023definition, else2023collisionless, shi2023loop}. Most importantly for the purpose of our work, an effective field theory (EFT) for FL was put forward in Ref.~\cite{Son2022} based on a novel bosonization method for the Fermi surface. It was shown that starting from any semiclassical phase space distribution $f^{\text{sc}}(\mathbf{p};t,\mathbf{x})$ that represents a sharp Fermi surface, all other possible Fermi liquid distributions can be obtained by a phase space volume-preserving transformation parametrized by a bosonic field $\phi(\theta;t,\mathbf{x})$ that depends on spacetime coordinates and the position on the Fermi surface $\theta$. An effective action that reproduces known results at quadratic order and systematically includes non-linear corrections was further proposed for these bosonic degrees of freedom. Constructing an EFT for FL that captures non-perturbative effects and allows one to calculate all desired observables at any given order is a significant step towards an analytical description of NFL. However, it remains \emph{a priori} unclear if the ideas put forward by Ref.~\cite{Son2022} can be adapted to describe the infrared behavior of NFL, whose physics is fundamentally different from FL. 

Inspired by this framework,  we propose a bosonized effective field theory for NFL. It is first shown that, in light of the quantum Boltzmann equation (QBE), one should use the frequency-dependent generalized distribution $f(\omega,\theta;t,\mathbf{x})$ rather than the semiclassical phase space distribution $f^{\text{sc}}(\mathbf{p};t,\mathbf{x})$ for the dynamics of NFL. Here, $\theta$ represents the angular coordinate along the FS in momentum space and $\omega$ the frequency. Indeed, it is highlighted that the usual distribution function over phase space does not satisfy a closed equation of motion in the absence of quasiparticles. In contrast, the generalized distribution $f(\omega,\theta;t,\mathbf{x})$ importantly captures the nontrivial frequency dependence inherent to any description of NFL and does follow a closed equation of motion even in the absence of clearly defined excitations. We then argue that physical fluctuations of these generalized distributions can be parametrized by a bosonic field and introduce an effective action. We demonstrate that this action yields the usual collisionless QBE as its equation of motion. Expanding the action to quadratic order shows that our effective theory reproduces known results for the case of a two-dimensional spherical Fermi surface. The action shows forward scattering cancellation, which can usually be attributed to a Ward identity~\cite{kim1994gauge, Kim1995}. By performing an angular momentum decomposition of the bosonic field, it is further proven that the small angular momentum modes display the usual Fermi liquid-like behavior. In contrast, the large angular momentum modes (or rough fluctuations) have a singular behavior characteristic of NFL --- a known result typically obtained through a detailed analysis of the QBE~\cite{Kim1995, Mandal2022}. This effective action thus seemingly captures highly nontrivial known results even if restricted to quadratic order. It should be emphasized that even though we focus in this article on the quadratic action for a spherical Fermi surface in two spatial dimensions, our proposed effective action applies to other nontrivial Fermi surfaces in any dimension and, most importantly, systematically captures non-linear corrections that would be intractable with other means. By offering a comprehensive framework to describe the physics of NFL analytically, our work should serve as an essential cornerstone towards a more complete understanding of strange metals and related problems. 

The rest of the paper is organized as follows: Section~\ref{sec:QBE_GDF} briefly reviews the derivation of the quantum Boltzmann equation and introduces the generalized distribution function. In Section~\ref{sec:effective_action}, we derive the effective action before expanding it to the quadratic order in terms of the bosonic field and computing response functions at first order in the self-energy expansion in Section~\ref{sec:quadatic}. We then comment on how to generalize the quadratic action using the full functional expansion of the self-energy in Section~\ref{sec:general}. In Section~\ref{sec:discussion}, we finally end by discussing our work's implications and possible future directions.



%%%%%%%%%%%%%%%%%%%%%%%%%%%%%%%%%%%%%%%%%%%%%%%%%%%%%%%%%%%%%%%%%%%%%%%%%%%
\section{Quantum Boltzmann equation and generalized distribution function}\label{sec:QBE_GDF}

Much of the current understanding of FL and NFL relies on using the quantum Boltzmann equation. We here introduce the usual derivation of these equations of motion before discussing how they can naturally emerge from an effective field theory in the following sections. 

The QBE describes the evolution of the fermionic distribution function over time and can therefore be thought of as a coarse-grained version of the equation of motion for the non-equilibrium one-particle lesser Green's function
\begin{equation}
    G^{<}\left(\mathbf{x}_1, t_1; \mathbf{x}_2, t_2\right) = i\left\langle\psi^{\dagger}\left(\mathbf{x}_2, t_2\right) \psi\left(\mathbf{x}_1, t_1 \right)\right\rangle. 
\end{equation}
In the general non-equilibrium case, the Green's function does not have spatial or temporal translation symmetry. It depends on both absolute (i.e., $\mathbf{x}=(\mathbf{x}_1+\mathbf{x}_2)/2$ and $t=(t_1+t_2)/2$) and relative coordinates (i.e., $\mathbf{x}_{\mathrm{rel}}=\mathbf{x}_{2} - \mathbf{x}_1$ and $t_{\mathrm{rel}}=t_{2}-t_{1}$). The equation of motion in the collisionless limit for the Fourier transform of the lesser Green's function with respect to the relative coordinates $G^{<}(\omega,\mathbf{p};t,\mathbf{x})$ is (see Appendix~\ref{app:Non-equilibirum_field_theory})~\cite{kita2010introduction, rammer1986quantum, rammer2011quantum}
\begin{equation}
\label{eq:Dyson_equation_general}
    [ \omega - \xi_{\mathbf{p}} -\Re\Sigma^{R}, G^{<}] - [\Sigma^{<}, \Re G^{R}] = 0,
\end{equation}
where we have introduced the free-particle spectrum $\xi_{\mathbf{p}}=\epsilon(\mathbf{p})-\mu$, the retarded Green's function $G^{R}$, its corresponding self-energy $\Sigma^{R}$, and the generalized Poisson bracket 
\begin{align} \label{eq:generalized_Poisson_bracket}
[A,B]=\frac{\partial A}{\partial\omega}\frac{\partial B}{\partial t}-\frac{\partial A}{\partial t}\frac{\partial B}{\partial\omega}-\frac{\partial A}{\partial p_{i}}\frac{\partial B}{\partial x_{i}}+\frac{\partial A}{\partial x_{i}}\frac{\partial B}{\partial p_{i}}
\end{align}
with the repeated indices summed over. This generalization contains additional derivatives for the relative frequency $\omega$ and absolute time $t$ in contrast to the usual Poisson bracket. In equilibrium, the lesser Green's function takes the form~\cite{bruus2004many, negele2018quantum, mahan2000many}
\begin{equation}
    G^{<}(\omega,\mathbf{p}; t,\mathbf{x}) = i A(\omega,\mathbf{p}; t,\mathbf{x}) f_{0}(\omega), 
\end{equation}
where $f_{0}(\omega)=\left( e^{\omega/T} + 1 \right)^{-1}$ is the Fermi-Dirac distribution and 
\begin{equation}
    A(\mathbf{p}, \omega)=\frac{-2 \operatorname{Im} \Sigma^R(\omega, \mathbf{p})}{\left(\omega-\xi_{\mathbf{p}}-\operatorname{Re} \Sigma^R(\omega, \mathbf{p})\right)^2+\left(\operatorname{Im} \Sigma^R(\omega, \mathbf{p})\right)^2}.
\end{equation}
is the spectral function. 

For Fermi liquids, quasiparticles remain well-defined even in the presence of strong interactions such that $\operatorname{Im} \Sigma^R\ll \omega$ at a sufficiently low energy scale. In this physically relevant low-energy limit, we can neglect the self-energy and approximate the quasiparticles to be infinitely long-lived~\cite{bruus2004many, negele2018quantum, mahan2000many}
\begin{align}
    A(\omega, \mathbf{p}) &\approx 2 \pi \delta\left(\omega-\xi_{\mathbf{p}}  \right). 
\end{align}
In the presence of quasiparticles, we can thus integrate out the frequency dependence of Eq.~\eqref{eq:Dyson_equation_general} to obtain a closed equation of motion  for the \emph{semiclassical distribution} function
\begin{equation}
    f^{\text{sc}}(\mathbf{p}; t, \mathbf{x}) = -i \int \frac{d \omega}{2 \pi} G^{<}(\omega,\mathbf{p}; t,\mathbf{x})
\end{equation}
which is the usual collisionless Boltzmann equation
\begin{equation} \label{eq:QBE_for_FL}
    \left[ \omega - \xi_{\mathbf{p}}, f^{\text{sc}} \right] = \partial_{t}f^{\text{sc}} +(\nabla_{\mathbf{p}}\epsilon(\mathbf{p}))\cdot(\nabla_{\mathbf{x}} f^{\text{sc}}) = 0. 
\end{equation}
A superscript ``sc'' emphasizes that this is analogous to the usual semiclassical phase space distribution.

Interactions are so important for NFL that the concept of quasiparticles breaks down since $\operatorname{Im} \Sigma^R$ scales as fast or faster than $\omega$ at low energy. Furthermore, the corresponding $\text{Re}\Sigma^{R}$ leads to a divergent effective mass of the quasiparticles. The spectral function is no longer strongly peaked, implying that, even near equilibrium, the frequency-independent distribution function $f^{\text{sc}}(\mathbf{p}; t, \mathbf{x})$ does not satisfy a closed equation of motion anymore. An accurate description of the system thus requires proper consideration of the frequency dependence of the self-energy. In fact, the frequency dependence of the self-energy often dominates its contribution to Eq.~\eqref{eq:Dyson_equation_general} such that we can assume at leading order that $\Sigma^{R}$ only depends on $\omega$ and is independent of the magnitude of momentum measured from the Fermi surface. In this limit, the spectral function is a well-peaked function of $\xi_{\mathbf{p}}$ at the Fermi surface for small frequency~\cite{Kadanoff1964,Kim1995}. Even in the absence of well-defined quasiparticles, an equation of motion can be obtained by integrating Eq.~\eqref{eq:Dyson_equation_general} over $\xi_{\mathbf{p}}$ for the \emph{generalized distribution} function
\begin{equation}
    f(\omega, \theta ; t, \mathbf{x})=-i \int \frac{d \xi_\mathbf{p}}{2 \pi} G^{<}(\omega,\mathbf{p}; t,\mathbf{x}),
\end{equation}
where $\theta$ parametrizes the direction of $\mathbf{p}$. The equation of motion for the generalized distribution function is the collisionless QBE  
\begin{equation}\label{eq:QBE_for_NFL}
    \left[\omega-\xi_{\mathbf{p}}-\Sigma^{R}[f],f \right] = 0,
\end{equation}
where we used $\int d\xi_{\mathbf{p}} \Re
G^{R}=0$ (see Appendix~\ref{app:useful}) and consider the linearized dispersion \begin{equation} \label{eq:linearized_dispersion}
    \xi_{\mathbf{p}} = \epsilon(\mathbf{p})-\mu \approx v_{F}(\theta) p_{\parallel}. 
\end{equation}
Here, $p_\parallel$ is the momentum component parallel to $\nabla_{\mathbf{p}}\epsilon$ (i.e., perpendicular to the Fermi surface), and the Fermi velocity is $v_{F}(\theta) \mathbf{n}_{\theta}=\left.\nabla_{\mathbf{p}} \epsilon \right|_{\theta, \xi_p=0}$. Notice that $\nabla_{\mathbf{p}}f\approx p_{F}^{-1}\mathbf{s}_{\theta}\partial_{\theta}f$ since $f$ only depends on the angle $\theta$ or the momentum transfer parallel to the Fermi surface. From this point onward, we will consider the self-energy to be real and shall accordingly not mention that we are taking the real part to lighten the notation as in Eq.~\eqref{eq:QBE_for_NFL}. The generalized distribution can be thought of as describing the variation of the local chemical potential in momentum space. Linearizing this equation of motion, one can derive certain properties of NFL. However, a much deeper and more systematic understanding of NFL could be gained from an effective action with Eq.~\eqref{eq:QBE_for_NFL} as its saddle point solution. In the next section, we propose such an EFT.


A central object in the equation of motion for NFL is the self-energy. $\Sigma^{R}$ is a functional of the generalized distribution $f$ and can be expanded as
\begin{widetext}
\begin{align}
\Sigma^{R}[f]={}&\int_{\omega_{1}\theta_{1}} \frac{\delta\Sigma^{R}}{\delta f(\omega_{1},\theta_{1})} f(\omega_{1},\theta_{1})+\frac{1}{2!}\int_{\omega_{1}\theta_{1}}\int_{\omega_{2}\theta_{2}}\frac{\delta^{2}\Sigma^{R}}{\delta f(\omega_{1},\theta_{1})\delta f(\omega_{2},\theta_{2})} f(\omega_{1},\theta_{1}) f(\omega_{2},\theta_{2})+\ldots\notag\\
={}&\int_{\omega_{1}\theta_{1}}D_{1}(\omega_{1},\theta_{1};\omega,\theta) f(\omega_{1},\theta_{1})+\int_{\omega_{1}\theta_{1}}\int_{\omega_{2}\theta_{2}}D_{2}(\omega_{1},\theta_{1},\omega_{2},\theta_{2};\omega,\theta) f(\omega_{1},\theta_{1}) f(\omega_{2},\theta_{2})+\ldots=\sum_{n=1}^{\infty} \Sigma^{R(n)}[f]\label{eq:self_expand},
\end{align}
\end{widetext}
where we have used the notation $f(\omega_{n},\theta_{n})\equiv f(\omega_{n},\theta_{n};t,x)$ for simplicity. The coefficients $D_{n}=\frac{1}{n!}\frac{\delta^{n}\Sigma^{R}}{\delta f(\omega_{1},\theta_{1})\ldots \delta f(\omega_{n},\theta_{n})}$ are the generalized Landau-interaction functions. We assume that these generalized coefficients do not depend on position or time, are invariant under the change of variables 
\begin{equation} \label{eq:Change_of_variables_Landau_parameters}
    D_{n}(\ldots, \omega_j,\theta_j,\ldots;\omega,\theta) = D_{n}(\ldots,\omega,\theta, \ldots;\omega_j,\theta_j),
\end{equation} 
and only depend on the frequency and angle differences $D_{n}(\omega_{1},\theta_{1},\cdots;\omega,\theta)=D_{n}(\omega_{1}-\omega,\theta_{1}-\theta;\cdots;\omega_{n}-\omega,\theta_{n}-\theta)$.
It is important to note the frequency dependence of $D_{n}$ in place of the usual static coefficients used in Landau's Fermi liquid theory. For the momentum integral, we have only kept the main contribution coming from the integration close to the Fermi surface $\int_{\omega\theta}=\int\frac{d\omega d^{d-1}\theta}{(2\pi)^{d}}N(\mu, \theta)$ where $N(\epsilon =\mu,\theta)$ is the density of states near the Fermi surface at an angle $\theta$. 
The $n$th-order term in the self-energy expansion $\Sigma^{R(n)}$ is
\begin{widetext}
\begin{align}
\Sigma&^{R(n)}[f]=%\notag\\={}&
\int_{\omega_{1}\theta_{1}}\ldots\int_{\omega_{n}\theta_{n}}D_{n}(\omega_{1}-\omega,\theta_{1}-\theta;\ldots;\omega_{n}-\omega,\theta_{n}-\theta)%\notag\\&\quad\times
\prod_{i=1}^{n}f(\omega_{i},\theta_{i}).
\end{align}
\end{widetext}


%%%%%%%%%%%%%%%%%%%%%%%%%%%%%%%%%%%%%%%%%%%%%%%%%%%%%%%%%%%%%%%%%%%%
\section{Effective action via generalized canonical transformation}\label{sec:effective_action}

The equations of motion described in Eqs.~\eqref{eq:QBE_for_FL} and \eqref{eq:QBE_for_NFL} offer important insight into the properties of FL and NFL. They allow for the calculation of certain response functions at leading order for small momentum. However, to gain a deeper and more systematic understanding, it would be preferable to have a description of the low-energy behavior of these phases of matter in terms of an effective field theory. An effective action would provide a way to evaluate, in theory, all possible observables at any given order. For FL, such an effective action was recently proposed by Delacrétaz, Du, Mehta, and Son~\cite{Son2022} using the method of coadjoint orbits. Inspired by their work, we propose a bosonization approach suitable for NFL. 


%%%%%%%%%%%%%%%%%%%%%%%%%%%%%%%%%%%%%%%
\subsection{Generalized canonical transformation}

The quantum Boltzmann equations~\eqref{eq:QBE_for_FL} and~\eqref{eq:QBE_for_NFL} govern the evolution of any semiclassical $f^{\text{sc}}$ and generalized distribution functions $f$, respectively. It is widely accepted that both FL and NFL possess sharp Fermi surfaces at zero temperature, which persist even when the distribution undergoes time evolution under the collisionless QBE. Consequently, it suffices to consider distributions that describe these sharp Fermi surfaces to capture the fluctuations of FL and NFL. In the case of Fermi liquids, this corresponds to a distribution that exhibits a jump at the Fermi momentum from the occupied to the unoccupied side. For NFL, generalized distributions with Fermi surfaces are associated with Green's functions that manifest a singularity or kink at the Fermi momentum in their momentum occupation $n(\mathbf{p})$. By starting from any representative $f_0$ of this class of distribution functions, all other sharp distributions $f$ can be parametrized by a transformation $U$. 


For instance, a natural choice of representative generalized distribution is the equilibrium Fermi-Dirac distribution $f_{0}(\omega)=(e^{\omega/T}+1)^{-1}$ which reduces to the Heaviside step function $f_{0}(\omega)=\Theta(-\omega)$ at zero temperature. This system has a free-particle spectrum $\xi_{\mathbf{p}}$ and associated Fermi surface defined by the equation $\xi_{\mathbf{p}}=0$. Any other generalized distribution relevant for NFL at zero-temperature can then be obtained by acting with a transformation $U=\exp\left( -\phi\right)$ as
\begin{align} \label{eq:canonical_transformatioon_quasiclassical_distribution}
    f(\omega,\theta;t,\mathbf{x})=&f_{0}-[\phi,f_{0}]+\frac{1}{2!}[\phi,[\phi,f_{0}]]+\ldots \notag\\
    =& U f_{0} U^{-1},
\end{align}
where the generalized Poisson bracket introduced in Eq.~\eqref{eq:generalized_Poisson_bracket} has to be used since the distribution $f$ depends on frequency, time, position and momentum. One should think of $U$ not as the exponential of a function but rather as the exponential map of a Lie algebra element $\phi$, which is isomorphic to a smooth function that goes to zero at infinity $\phi\cong\phi(\omega, \theta; t, \mathbf{x})\in C^{\infty}(\mathbb{R}^{2d+1})$~\cite{Son2022}. Eq.~\eqref{eq:canonical_transformatioon_quasiclassical_distribution} then naturally follows from the Baker–Campbell–Hausdorff formula~\cite{hall2013lie, miller1973symmetry}. 

The deviation of the generalized distribution function from the equilibrium $\delta f=f-f_{0}$ can be written in terms of $\phi$ as $\delta f=\sum_{n=1}\Delta_{\phi}^{(n)}f$ where
\begin{align}\label{eq:expansion_qc_distribution_phi}
\Delta^{(n)}_{\phi}f={}&\frac{(-1)^{n}}{n!}[\underbrace{\phi,[\phi,\ldots,[\phi}_{n},f_{0}]\ldots]].
\end{align}
Given its functional dependence on the distribution function (see Eq.~\eqref{eq:First_order_sefl_energy}), the self-energy will also be displaced from its equilibrium configuration $\Sigma^{R}_0\equiv\Sigma^{R}[f_0]$ by 
\begin{equation}\label{eq:expansion_self-energy_distribution_phi}
    \delta\Sigma^{R}[f]= \Sigma^{R}[f]- \Sigma^{R}_0 =\sum_{n=1}^{\infty}\sum_{m=0}^{\infty}\Delta^{(m)}_{\phi}\Sigma^{R(n)},    
\end{equation}
where $\Delta_{\phi}^{(m)}\Sigma^{R(n)}$ is a shorthand notation for the sum of all terms that come from the $n$th order in the functional expansion of the self-energy $\Sigma^{R(n)}$ and include $m$ functions $\phi$.

%%%%%%%%%%%%%%%%%%%%%%%%%%%%%%%%%%%%%%%
\subsection{Effective action for Non-Fermi liquids}

We propose the following effective action to describe the low-energy physics of non-Fermi liquids
\begin{align}
S={}&\int_{t\omega} \braket{f,\omega-\xi_{\mathbf{p}}-\sum_{n=1}\tfrac{1}{n+1}\Sigma^{R(n)}[f]}\notag\\
={}&\int_{t\omega} \braket{f_{0},U^{-1}(\omega-\xi_{\mathbf{p}}-\sum_{n=1}\tfrac{1}{n+1}\Sigma^{R(n)}[f])U},\label{eq:action}
\end{align}
where $\int_{t\omega}=\int dtd\omega/2\pi$,  $\braket{A,B}$ is the inner product defined by
\begin{align}
    \braket{A,B}=\mathcal{C}\int \frac{d^{d}xd^{d}p}{(2\pi)^{d}}AB
\end{align}
and $\mathcal{C}$ is a normalization constant that reduces to $\mathcal{C}=N(\mu)/(\int\frac{d\epsilon}{2\pi}N(\epsilon))^{-1}$ in the isotropic limit.  In the action, $\xi_{\mathbf{p}}$ should be thought of as the linearized dispersion at the Fermi surface, just like in Eq.~\eqref{eq:linearized_dispersion}.
Eq.~\eqref{eq:action} is thus an effective field theory for the bosonic degrees of freedom $\phi$ that parametrize fluctuations of the generalized distribution.

To verify that the above effective action properly captures the dynamics of non-Fermi liquids, let us obtain the equation of motion $\delta S/\delta \phi=0$. To do so we vary the field $\phi$ by an infinitesimal amount $\alpha$ (i.e., $U\rightarrow e^{-\alpha}U$). The variations of the action at linear order in $\alpha$ are (see Appendix~\ref{app:general_action} for details)
\begin{align}
\delta[U^{-1}BU]={}&U^{-1}[\alpha,B]U,\label{eq:varA}\\
\delta[U^{-1}\Sigma^{R(n)}[f]U]={}&U^{-1}[\alpha,\Sigma^{R(n)}]U \notag\\
& - n\int_{\omega_{1} \theta_{1}}\cdots\int_{\omega_{n} \theta_{n}}  D_{n}[\alpha_{n},f_{n}]_{n}\prod_{i=1}^{n-1}f_{i} ,\label{eq:varS}
\end{align}
where $B$ stands for any parameters in the action without implicit dependence on the generalized distribution function, $\alpha_{i}$ and $f_{i}$ depend on $\omega_{i}$ and $\theta_{i}$, and $[A,B]_{n}$ indicates that the derivatives in the generalized Poisson bracket are evaluated with respect to $\omega_{n}$ and $\mathbf{p}_{n}$. By using Eqs.~\eqref{eq:varA} and \eqref{eq:varS}, the first-order variation of the action $\delta S$ is given by
\begin{align}
\delta S
={}&\int_{t\omega}\braket{[\omega-\xi_{\mathbf{p}}-\Sigma^{R}[f],f],\alpha},
\end{align}
where the identity $\int_{t\omega}\braket{A,[B,C]}=\int_{t\omega}\braket{[A,B],C}$ was used. % (see Appendix~\ref{app:useful} for details). 
As a result, the equation of motion is
\begin{align}
[\omega-\xi_{\mathbf{p}}-\Sigma^{R}[f],f]=0,\label{eq:QBE}
\end{align}
which is none other than the collisionless QBE for NFL introduced in Eq.~\eqref{eq:QBE_for_NFL}.

At first glance, the factor on the right of the inner product in the action~\eqref{eq:action} might look similar to the inverse Green's function. However, there is an extra factor of $(n+1)^{-1}$ in front of the $n$th order in $f$ of the self-energy expansion. The factor is required to reproduce the quantum Boltzmann equation from the action~\eqref{eq:action} since the functional dependence of the self-energy on the distribution function leads to extra contributions (see  Eq.~\eqref{eq:varS}). 
Also, Eq.~\eqref{eq:action} can be rewritten as
\begin{align}\label{eq:action_luttinger-Ward}
    S={}&\int_{t\omega}\braket{f,\omega-\xi_{\mathbf{p}}}-\int_{tx}\Phi[f],
\end{align}
where we have introduced the functional
\begin{align}\label{eq:definition_generalized_luttinger-Ward}
\Phi[f]={}\sum_{n=1}^{\infty}\frac{1}{n+1}\int_{\omega\theta}\Sigma^{R(n)}[f] f.  
\end{align}
Interestingly, the functional derivative of $\Phi[f]$ with respect to the distribution satisfies $\delta\Phi/\delta f=\Sigma^{R}$, which is completely analogous to the conventional Luttinger-Ward functional~\cite{luttinger1960ground, baym1961conservation, kotliar2006electronic}. This analogy can even be pushed further by noting that the Feynman diagrams produced by this term are in correspondence with some two-particle irreducible bubble diagrams (i.e., ``skeleton'' diagrams~\cite{potthoff2004non, kozik2015nonexistence}) of the initial fermion theory at equilibrium (see Appendix~\ref{app:LW} for details). Considering the analogy, we shall refer to $\Phi$ as the generalized Luttinger-Ward functional of our effective field theory. 

In sum, the proposed action~\eqref{eq:action} produces the collisionless quantum Boltzmann equation.  
It does not assume the existence of well-defined quasiparticles and captures the non-trivial frequency dependence of NFL through the self-energy and its use of the generalized instead of the semiclassical distribution function. It should be underscored that this action contains terms at all orders in $\phi$.


%%%%%%%%%%%%%%%%%%%%%%%%%%%%%%%%%%%%%%%%%%%%%%%%%%%%%%%%%%%%%%%%%%%%%%%%%%%
\section{Quadratic theory at one loop order}\label{sec:quadatic}

The action proposed in the last section describes a non-linear effective theory for the $\phi$ field that parametrizes fluctuations of the generalized distribution around its equilibrium position. In this section, we discuss how to expand the effective action in terms of $\phi$ and use it to compute observables. 
In particular, we will study the quadratic action in $\phi$ assuming a spherical Fermi surface in two dimensions to compare with existing results. 
In addition, we will consider the first-order expansion in $f$ for the self-energy (i.e., one-loop approximation)
\begin{align}\label{eq:First_order_sefl_energy}
    \Sigma^{R}[f]={}&\int_{\omega'\theta'} D_{1}(\omega',\omega;\theta',\theta)f(\omega',\theta') = \Sigma^{R(1)}[f].
\end{align}
It is further assumed that the first order generalized coefficient only depends on the frequency and angular separations $D_{1}(\omega',\theta';\omega,\theta)=D_{1}(\omega'-\omega,\theta'-\theta)$ and is invariant under the change of variables $D_{1}(\omega'-\omega,\theta'-\theta)=D_{1}(\omega-\omega',\theta'-\theta)=D_{1}(\omega'-\omega,\theta-\theta')\equiv D_1$. Such assumptions are accurate in many physical situations of interest like a Fermi surface in two dimensions interacting with a transverse gauge boson~\cite{kim1994gauge,Kim1995} or a critical Ising-nematic order parameter~\cite{Mandal2022}. In these cases, the one-loop random phase approximation (RPA) boson propagator corresponds to $D_{1}$ and is approximately given by $D_{1}(\omega'-\omega,\mathbf{p}'-\mathbf{p})\approx D_{1}(\omega'-\omega,k_{F}|\theta'-\theta|)$ since the main contribution comes from near the Fermi surface.
We emphasize that the action can be used for more general Fermi surfaces in any dimension and systematically captures higher-order contributions in $\phi$ for the effective action and in $f$ for the self-energy.
The general case, which includes the higher-order contributions in $f$ for the self-energy, will be discussed in the next section.


%%%%%%%%%%%%%%%%%%%%%%%%%%%
\subsection{Derivation of the Gaussian action}\label{subsec:quadatic_gaussian}

Using the generalized canonical transformations for the generalized distribution~\eqref{eq:canonical_transformatioon_quasiclassical_distribution} as well as the definitions given in Eqs.~\eqref{eq:expansion_qc_distribution_phi} and~\eqref{eq:expansion_self-energy_distribution_phi} for the distribution and self-energy displacement, the action within the one-loop approximation can be expanded as
\begin{widetext}
\begin{align}
S={}&\int_{t\omega}\braket{f,\omega-\xi_{\mathbf{p}}-\tfrac{1}{2}\Sigma^{R}[f]}\notag\\
={}&\int_{t\omega} \braket{f_{0}+\Delta_{\phi}^{(1)}f+\Delta_{\phi}^{(2)}f+\ldots,\omega-\xi_{\mathbf{p}}-\tfrac{1}{2}(\Sigma^{R}[f_0]+\Delta_{\phi}^{(1)}\Sigma^{R}+\Delta_{\phi}^{(2)}\Sigma^{R}+\ldots)}\notag\\
={}&\int_{t\omega} \braket{f_{0}(\omega),\omega-\xi_{\mathbf{p}}-\tfrac{1}{2}\Sigma^{R}_{0}(\omega)}\notag\\
%%%%
&+\int_{t\omega} \braket{\Delta_{\phi}^{(1)}f,\omega-\xi_{\mathbf{p}}-\Sigma^{R}_{0}(\omega)}\notag\\
%%%%
&+\int_{t\omega}\braket{\Delta_{\phi}^{(2)}f,\omega-\xi_{\mathbf{p}}-\Sigma^{R}_{0}(\omega)}
-\int_{t\omega}\braket{\Delta_{\phi}^{(1)}f,\tfrac{1}{2}\Delta_{\phi}^{(1)}\Sigma^{R}}\notag\\
%%%%
&+\ldots,\label{eq:effective_expand}
\end{align}
\end{widetext}
where $\Sigma_{0}^{R}=\Sigma^{R(1)}[f_{0}]$, $\Delta_{\phi}^{(m)}\Sigma^{R}=\Delta_{\phi}^{(m)}\Sigma^{R(1)}$, and we used $\int_{t\omega}\braket{A,[B,C]}=\int_{t\omega}\braket{[A,B],C}$ and $\int_{t\omega}\braket{\Delta_{\phi}^{(m)}f,\Delta_{\phi}^{(n)}\Sigma^{R}}
={}\int_{t\omega}\braket{\Delta_{\phi}^{(n)}f,\Delta_{\phi}^{(m)}\Sigma^{R}}$  (see Appendix~\ref{app:useful} for details). The first and second lines of the last equality correspond to the $\phi^{0}$ and $\phi^{1}$ orders, respectively. The action at the $\phi^{0}$ order is a constant, and the $\phi^{1}$ order contribution gives a vanishing boundary term as it corresponds to a total time derivative. The third line is the quadratic action $S^{(2)}$, whereas the ellipsis in the last line stand for higher-order corrections that we discuss in more detail in Appendix~\ref{app:higher}. 

The Gaussian action for $\phi$, $S^{(2)}$, can be split into two: the non-interacting $S^{(2)_{0}}$ and interacting $S^{(2)_{\text{int}}}$ parts
\begin{align}
S^{(2)}={}&S^{(2)_{0}}+S^{(2)_{\text{int}}},\\
%%%
S^{(2)_{0}}={}&-\frac{1}{2}\int_{tx\omega\theta}(\partial_{\omega}f_{0})\dot{\phi}[\dot{\phi}+v_{F}\mathbf{n}_{\theta}\cdot(\nabla_{\mathbf{x}}\phi)],\label{eq:cont_quad_bare}\\
S^{(2)_{\text{int}}}={}&\frac{1}{2}\int_{tx\omega\theta}(\partial_{\omega}f_{0})\dot{\phi}[(\partial_{\omega}\Sigma^{R}_{0})\dot{\phi}-\Delta^{(1)}_{\phi}\Sigma^{R}].\label{eq:cont_quad_int}
\end{align}
We note that the Fermi velocity and density of states are now isotropic since we are considering a spherically symmetric Fermi surface (i.e., $v_{F}=\epsilon'(p)|_{p=p_{F}}$, $\mathbf{n}_{\theta}=\mathbf{p}_{F}/p_{F}$, and $N(\epsilon,\mathbf{0})= \int d\theta/(2\pi) N(\epsilon,\theta)\equiv N(\epsilon)$) 

We can check the equation of motion for the quadratic Gaussian action. The linear order variation for $\phi\rightarrow\phi+\delta\phi$ is
\begin{align}
    \delta S^{(2)}={}&\int_{t\omega}\braket{[\omega-\xi_{\mathbf{p}}-\Sigma^{R}_{0},\Delta_{\phi}^{(1)}f]-[\Delta_{\phi}^{(1)}\Sigma^{R},f_{0}],\delta\phi},
\end{align}
which leads to the equation of motion 
\begin{align}
    [\omega-\epsilon-\Sigma^{R}_{0},\Delta_{\phi}^{(1)}f]-[\Delta_{\phi}^{(1)}\Sigma^{R},f_{0}]=0.
\end{align}
This is simply the QBE for NFL expanded at the $\phi^{1}$ order.


In terms of the Fourier transformed field $\hat{\phi}(\omega,\theta;\Omega,\mathbf{q})=\int_{tx}e^{-i(\Omega t-\mathbf{q}\cdot\mathbf{x})}\phi(t,x;\omega,\theta)$, the two parts of the Gaussian actions are
\begin{widetext}
\begin{align}
    S^{(2)_{0}}={}&-\frac{1}{2}\int_{\Omega q\omega\theta}F_{0}(\omega,\Omega)(\Omega-v_{F}(\mathbf{n}_{\theta}\cdot\mathbf{q}))\hat{\phi}(\Omega,\mathbf{q};\omega,\theta)\hat{\phi}(-\Omega,-\mathbf{q};\omega,\theta),\label{eq:disc_quad_bare}\\
   S^{(2)_\text{int}}
={}&\frac{1}{2}\int_{\Omega q\omega\theta}F_{0}(\omega,\Omega)(\Sigma^{R}_{0}(\omega+\Omega)-\Sigma^{R}_{0}(\omega))\hat{\phi}(\Omega,\mathbf{q};\omega,\theta)\hat{\phi}(-\Omega,-\mathbf{q};\omega,\theta)\notag\\
&-\frac{1}{2}\int_{\Omega q\omega\theta}\int_{\omega'\theta'}F_{0}(\omega,\Omega)F_{0}(\omega',\Omega)D_{1}(\omega'-\omega,\theta'-\theta)\hat{\phi}(\Omega,\mathbf{q};\omega,\theta)\hat{\phi}(-\Omega,-\mathbf{q};\omega',\theta')\label{eq:disc_quad_int}
\end{align}
where we discretized the frequency derivative ($(\partial_{\omega}f_{0})\Omega\approx f_{0}(\omega+\Omega)-f_{0}(\omega)$ for small $\Omega$) for regularization (see Appendix~\ref{app:quadratic} for details) and used a shorthand notation $F_{0}(\omega,\Omega)\equiv f_{0}(\omega+\Omega)-f_{0}(\omega)$. Using Eq.~\eqref{eq:First_order_sefl_energy} for the equilibrium self-energy, the interacting part $S^{(2)_{\text{int}}}$ can be rewritten as
\begin{align}
S^{(2)_{\text{int}}}={}&\frac{1}{2}\int_{\Omega q\omega\theta}\int_{\omega'\theta'}F_{0}(\omega,\Omega)F_{0}(\omega',\Omega)D_{1}(\omega'-\omega,\theta'-\theta)\hat{\phi}(\Omega,\mathbf{q};\omega,\theta)
[\hat{\phi}(-\Omega,-\mathbf{q};\omega,\theta)-\hat{\phi}(-\Omega,-\mathbf{q};\omega',\theta')].
\end{align}
\end{widetext}
By inspection, one can notice that the interacting part of the quadratic action vanishes for forward scattering (i.e., near $\omega=\omega'$ and $\theta=\theta'$). In such a regime, the system thus exhibits Fermi liquid behavior. In the study of a Fermi surface coupled to a transverse gauge boson, it has already been pointed out that forward scattering cancellation arises from a Ward identity~\cite{kim1994gauge} and can be observed in the QBE~\cite{Kim1995}. The simplicity by which this non-trivial result can be obtained with our formalism graciously demonstrates its usefulness. In the following subsections, we shall study this quadratic action in detail to see how its predictions compare with existing results from the literature. 


%%%%%%%%%%%%%%%%%%%%%%%%%%%
\subsection{Two-point density correlation function in non-interacting limit}

\begin{figure}[ht]
    \centering
    \subfigure[]{\includegraphics[scale=0.9]{fig1a.pdf}\label{fig:Pi0}}\hfill
    \subfigure[]{\includegraphics[scale=0.9]{fig1b.pdf}\label{fig:Pi1}}\hfill
    \subfigure[]{\includegraphics[scale=0.9]{fig1c.pdf}\label{fig:Pi2}}\hfill
    \subfigure[]{\includegraphics[scale=0.9]{fig1d.pdf}\label{fig:Pi3}}
    \caption{Feynmann diagrams for two-point density correlation function}
    \label{fig:self}
\end{figure}

Using the effective action, we can compute observables such as the density correlation function of the fermions. Let us consider the propagator of the bosonic field $\phi$. Since we now consider the action up to quadratic order in $\phi$, we can find, in principle, the exact propagator of the Gaussian action by taking its inverse. However, because the interaction part $S^{(2)_{\text{int}}}$ of the action corresponds to a dense infinite dimensional matrix in terms of $\omega$ and $\theta$, obtaining the exact propagator is problematic. For that reason, we compute correlation functions by considering the interaction part perturbatively in reference to the non-interacting action. The bare propagator associated with the non-interacting action~\eqref{eq:disc_quad_bare} is 
\begin{widetext}
\begin{align}
G_{0}(\Omega,\mathbf{q};\omega,\theta;\omega',\theta')={}&-\frac{(2\pi)^{2}}{N(\mu)}\frac{\delta(\omega'-\omega)\delta(\theta'-\theta)}{(f_{0}(\omega+\Omega)-f_{0}(\omega))(\Omega-v_{F}(\mathbf{n}_{\theta}\cdot\mathbf{q}))}.\label{eq:non-int_prop}
\end{align}
\end{widetext}


The density operator must be expressed in terms of $\phi$ to compute the density correlation function.  Similar to the generalized distribution function, the deviation of the density operator $\delta\rho$ can be written in the leading order as
\begin{align}
    \delta\rho(t,\mathbf{x})={}&\int_{\omega\theta}\delta f(t,\mathbf{x};\omega,\theta)%\notag\\
    \approx{}-\int_{\omega\theta}[\phi,f_{0}]=\int_{\omega\theta}(\partial_{\omega}f_{0})\dot{\phi},
\end{align}
and its Fourier transform is
\begin{align}
    \delta\hat{\rho}(\Omega,\mathbf{q})\approx i\int_{\omega\theta}(\partial_{\omega}f_{0})\Omega\hat{\phi}
    \approx i\int_{\omega\theta}(f_{0}(\omega+\Omega)-f_{0}(\omega))\hat{\phi}. \label{eq:den_op}
\end{align}
By using Eqs.~\eqref{eq:non-int_prop} and \eqref{eq:den_op}, the two-point density correlation function is
\begin{align}
\braket{\delta\hat{\rho}\,\delta\hat{\rho}}_{0}(\Omega,\mathbf{q})={}&\int_{\omega\theta}\frac{f_{0}(\omega+\Omega)-f_{0}(\omega)}{\Omega-v_{F}q\cos\theta},\label{eq:non-int_den_correl}
\end{align}
where $\cos\theta=\mathbf{n}_{\theta}\cdot\hat{q}$. We can check that this is consistent with the particle-hole bubble contribution (see Fig.~\ref{fig:Pi0}) in the equilibrium theory of a Fermi surface with the integration over $\xi_{\mathbf{p}}$ performed first. The one-loop particle-hole bubble contribution for small momentum transfer $\mathbf{q}$ in imaginary time is  
\begin{align}
\Pi_{0}&(i\nu,\mathbf{q})\notag\\
={}&-\int\frac{d\nu'}{2\pi}\frac{d^{2}p}{(2\pi)^{2}} \frac{1}{i\nu'+i\nu-\xi_{\mathbf{p}+\mathbf{q}}}\frac{1}{i\nu'-\xi_{\mathbf{p}}}\notag\\
\approx{}&-i\frac{N(\mu)}{2}\int\frac{d\nu' d\theta}{(2\pi)^{2}} \frac{\text{sgn}(\omega_{n}+\nu)-\text{sgn}(\omega_{n})}{i\nu-v_{F}q\cos\theta}.
\end{align}
After a Wick rotation, it can be written as
\begin{align}
    \Pi(\Omega,\mathbf{q})\approx{}&\int_{\omega\theta}\frac{f_{0}(\omega+\Omega)-f_{0}(\omega)}{\Omega-v_{F}q\cos\theta},
\end{align}
which is the same as Eq.~\eqref{eq:non-int_den_correl}. Some subtleties regarding the order of integration have to be carefully considered when evaluating these expressions. We discuss these in Appendix~\ref{app:int_order}. 
Therefore, our effective Gaussian action in the non-interacting limit $S^{(2)_0}$ leads to the usual characteristic Fermi liquid behavior.

%%%%%%%%%%%%%%%%%%%%%%%%%%%
\subsection{Two-point density correlation function in the perturbative interacting limit}

We now consider the perturbative correction to the density correlation function. From Eq.~\eqref{eq:disc_quad_int}, the bilinear interaction vertices are given by
\begin{widetext}
\begin{align}
V_{\text{diag}}(\Omega,\mathbf{q};\omega,\theta;\omega',\theta')={}&-\frac{(2\pi)^{2}}{N(\mu)}(f_{0}(\omega+\Omega)-f_{0}(\omega))(\Sigma^{R}_{0}(\omega+\Omega)-\Sigma^{R}_{0}(\omega))\delta(\omega'-\omega)\delta(\theta'-\theta),\\
V_{\text{off}}(\Omega,\mathbf{q};\omega,\theta;\omega',\theta')={}&(f_{0}(\omega+\Omega)-f_{0}(\omega))(f_{0}(\omega'+\Omega)-f_{0}(\omega'))D_{1}(\omega'-\omega,\theta'-\theta).
\end{align}
Then, the leading order contributions to the two-point correlation function from the diagonal and off-diagonal perturbations are respectively
\begin{align}
\braket{\delta\hat{\rho}\,\delta\hat{\rho}}_{1a}(\Omega,\mathbf{q})={}&\int_{\omega\theta}\frac{\Sigma^{R}_{0}(\omega+\Omega)-\Sigma^{R}_{0}(\omega)}{(\Omega-v_{F}q\cos\theta)^{2}}(f_{0}(\omega+\Omega)-f_{0}(\omega)),\label{eq:perturb_a_EFT}\\
%%%%%%%
\braket{\delta\hat{\rho}\,\delta\hat{\rho}}_{1b}(\Omega,\mathbf{q})={}&-\int_{\omega_{1}\theta_{1}}\int_{\omega_{2}\theta_{2}}\frac{(f_{0}(\omega_{1}+\Omega)-f_{0}(\omega_{1}))(f_{0}(\omega_{2}+\Omega)-f_{0}(\omega_{2}))}{(\Omega-v_{F}q\cos\theta_{1})(\Omega-v_{F}q\cos\theta_{2})}D_{1}(\omega_{2}-\omega_{1},\theta_{2}-\theta_{1}).\label{eq:perturb_b_EFT}
\end{align}
Similar to the non-interacting limit, $\braket{\delta\rho\,\delta\rho}_{1a}$ and $\braket{\delta\rho\,\delta\rho}_{1b}$ correspond to the two-loop contributions in the equilibrium theory of a Fermi surface interacting with bosons through a Yukawa coupling depicted in Fig.~\ref{fig:Pi1} and~\ref{fig:Pi2}, and Fig.~\ref{fig:Pi3}, respectively (with the integration over $\xi_{\mathbf{p}}$ performed first). To verify that Eqs.~\eqref{eq:perturb_a_EFT} and \eqref{eq:perturb_b_EFT} derived from our proposed effective action are consistent with the results obtained in the usual perturbative expansion of the equilibrium theory, we compute them explicitly. The two-loop contribution in Fig.~\ref{fig:Pi1} and~\ref{fig:Pi2} from the equilibrium theory for small momentum transfer $\mathbf{q}$ is
\begin{align}
\Pi_{1a}(i\nu,\mathbf{q})={}&-\int\frac{d\nu'}{2\pi}\frac{d^{2}p}{(2\pi)^{2}} \frac{\Sigma^{R}_{0}(i\nu'+i\nu)}{(i\nu'+i\nu-\xi_{\mathbf{p}+\mathbf{q}})^{2}(i\nu'-\xi_{\mathbf{p}})}
-\int\frac{d\nu'}{2\pi}\frac{d^{2}p}{(2\pi)^{2}} \frac{\Sigma^{R}_{0}(i\nu')}{(i\nu'+i\nu-\xi_{\mathbf{p}+\mathbf{q}})(i\nu'-\xi_{\mathbf{p}})^{2}}\notag\\
\approx{}&-\frac{i}{2}\int\frac{d\nu' d\theta}{(2\pi)^{2}}\frac{(\text{sgn}(\nu'+\nu)-\text{sgn}(\nu'))}{(i\nu-v_{F}q\cos\theta)^{2}}(\Sigma^{R}_{0}(i\nu'+i\nu)-\Sigma^{R}_{0}(i\nu')),
\end{align}
and the contribution in Fig.~\ref{fig:Pi3} is
\begin{align}
\Pi_{1b}(i\nu,\mathbf{q})={}&-\int\frac{d\nu_{1}}{2\pi}\frac{d^{2}p_{1}}{(2\pi)^{2}}\int\frac{d\nu_{2}}{2\pi}\frac{d^{2}p_{2}}{(2\pi)^{2}}\frac{D_{1}(i\nu_{2}-i\nu_{1},\theta_{2}-\theta_{1})}{(i\nu_{1}+i\nu-\xi_{\mathbf{p}_{1}+\mathbf{q}})(i\nu_{1}-\xi_{\mathbf{p}_{1}})(i\nu_{2}+i\nu-\xi_{\mathbf{p}_{2}+\mathbf{q}})(i\nu_{2}-\xi_{\mathbf{p}_{2}})}\notag\\
\approx{}&\frac{1}{4}\int\frac{d\nu_{1}d\theta_{1}}{(2\pi)^{2}}\int\frac{d\nu_{2}d\theta_{2}}{(2\pi)^{2}}\frac{(\text{sgn}(\nu_{1}+\nu)-\text{sgn}(\nu_{1}))(\text{sgn}(\nu_{2}+\nu)-\text{sgn}(\nu_{2}))}{(i\nu-v_{F}q\cos\theta_{1})(i\nu-v_{F}q\cos\theta_{2})}D_{1}(i\nu_{2}-i\nu_{1},\theta_{2}-\theta_{1}),
\end{align}
where once again we use imaginary time (or Euclidean space-time).  
After the Wick rotation back to real frequencies, both contributions can be written as
\begin{align}   
%%%\\
\Pi_{1a}(\Omega,\mathbf{q})\approx{}&\int\frac{d\omega d\theta}{(2\pi)^{2}}\frac{(f_{0}(\omega+\Omega)-f_{0}(\omega))}{(\Omega-v_{F}q\cos\theta)^{2}}(\Sigma^{R}_{0}(\omega+\Omega)-\Sigma^{R}_{0}(\omega)),\\
\Pi_{1b}(\Omega,\mathbf{q})\approx{}&-\int\frac{d\omega_{1}d\theta_{1}}{(2\pi)^{2}}\int\frac{d\omega_{2}d\theta_{2}}{(2\pi)^{2}}\frac{(f_{0}(\omega_{1}+\Omega)-f_{0}(\omega_{1}))(f_{0}(\omega_{2}+\Omega)-f_{0}(\omega_{2}))}{(\Omega-v_{F}q\cos\theta_{1})(\Omega-v_{F}q\cos\theta_{2})}D_{1}(\omega_{2}-\omega_{1},\theta_{2}-\theta_{1}).
\end{align}
Thus, the contributions from Fig.~\ref{fig:Pi1}-\ref{fig:Pi3} are consistent with $\braket{\delta\hat{\rho}\,\delta\hat{\rho}}_{1a}$ and $\braket{\delta\hat{\rho}\,\delta\hat{\rho}}_{1b}$. Recall that $D_1$ used in the above equations is the real part of the boson propagator in real frequency \cite{Kadanoff1964,Kim1995}, as mentioned in Sec.~\ref{sec:QBE_GDF}.
We can rewrite the total first-order correction as 
\begin{align}
\braket{\delta\hat{\rho}\,\delta\hat{\rho}}_{1a}+\braket{\delta\hat{\rho}\,\delta\hat{\rho}}_{1b}={}&\int_{\omega_{1}\theta_{1}}\frac{(f_{0}(\omega_{1}+\Omega)-f_{0}(\omega))}{(\Omega-v_{F}q\cos\theta_{1})}\notag\\
&\hspace{1cm}\times \int_{\omega_{2}\theta_{2}}\left[\frac{f_{0}(\omega_{2}+\Omega)-f_{0}(\omega_{2})}{(\Omega-v_{F}q\cos\theta_{1})}-\frac{f_{0}(\omega_{2}+\Omega)-f_{0}(\omega_{2})}{(\Omega-v_{F}q\cos\theta_{2})}\right]D_{1}(\omega_{2}-\omega_{1},\theta_{2}-\theta_{1}).
\end{align}
\end{widetext}
So for $\theta_{2}=\theta_{1}$ (forward scattering) and small momentum transfer, this perturbative contribution vanishes, and the system exhibits the Fermi liquid-like behaviors, as discussed previously. Note that this does not happen for large momentum transfer.

%%%%%%%%%%%%%%%%%%%%%%%%%%%
\subsection{Angular momentum decomposition}\label{subsec:quadratic_angular}

Previous studies of the QBE for NFL in two dimensions showed that collective modes display very different behaviors~\cite {Kim1995}. In particular, for small angular momentum modes of the Fermi surface, there is a small-angle (forward-scattering) cancellation between the self-energy and the Landau interactions, which leads to the usual FL behavior. In comparison, large angular momentum modes show an anomalous singular behavior characteristic of NFL.

This subsection examines how this behavior transpires in our EFT by considering the angular momentum decomposition of the Gaussian action. Since we are working in two dimensions, we can introduce the following representation of $\hat{\phi}$,
\begin{align}
\hat{\phi}(\Omega,\mathbf{q};\omega,\theta)={}&\sum_{l}e^{i\theta l}\tilde{\phi}(\Omega,\mathbf{q};\omega,l),\label{eq:ang_decomp}
\end{align}
%%%%%%%%%%%%%%%%%%%%%%%%%%%%%%%%%%%%%%%
where $l$ is the angular momentum (i.e., the conjugate variable of $\theta$). Note that Eq.~\eqref{eq:ang_decomp} is valid for two dimensions, and one has to use spherical harmonics and generalizations thereof for the angular decomposition in three dimensions and higher. Then, the non-interacting part of the quadratic action in the angular momentum decomposition is given by 
\begin{widetext}
\begin{align}\label{eq:ang_decomp_quad_bare}
S^{(2)_{0}}
={}&-\frac{1}{2}\sum_{l}\int_{\Omega q\omega}F_{0}(\omega,\Omega)\tilde{\phi}(-\Omega,-\mathbf{q};\omega,l)(\Omega\tilde{\phi}(\Omega,\mathbf{q};\omega,l)-\frac{v_{F}q}{2}[\tilde{\phi}(\Omega,\mathbf{q},l+1)+\tilde{\phi}(\Omega,\mathbf{q};\omega,l-1)])
\end{align}
and the interacting part is
\begin{align}
S^{(2)_\text{int}}
={}&\frac{1}{2}\sum_{l}\int_{\Omega q}\int_{\omega\omega'}\int_{\theta}F_{0}(\omega,\Omega)F_{0}(\omega',\Omega)D_{1}(\omega'-\omega,\theta)\tilde{\phi}(-\Omega,-\mathbf{q};\omega,l)[\tilde{\phi}(\Omega,\mathbf{q};\omega,l)-\cos(l\theta)\tilde{\phi}(\Omega,\mathbf{q};\omega',l)].
%%%%
\end{align}
\end{widetext}
%%%%%%%%%%%%%%%%%%%%%%%%%%%%%%%%%%%%%%%
If we consider the scaling dimension of the free part of the action $S^{(2)_{0}}$ as zero, we can obtain the scaling dimension of the interacting part $S^{(2)_{\text{int}}}$.
We will attach the coupling constant $g$ to $S^{(2)_{\text{int}}}$ and examine how it scales under the leading order renormalization group (RG). For concreteness, let us consider a case similar to Refs.~\cite{kim1994gauge, Kim1995} where a two-dimensional fermionic system interacts with transverse gauge bosons whose propagator is of the form 
\begin{align}
    \mathcal{D}(\mathbf{q}, \omega)=\frac{1}{-i \gamma \frac{\omega}{|\mathbf{q}|}+\chi q^\eta},
\end{align}
where $\gamma$ and $\chi$ are constants whereas $1<\eta\leq2$. Assuming $[\omega]=z$, we have $[\theta]=z/(\eta+1)$ and $[D_{1}]=-\eta[\theta]$. 
Then, we can define the critical angular momentum $l_{c}\propto k_{F}^{-1}\Omega^{1/(1+\eta)}$, and we can classify the behaviors depending on small $(l<l_{c})$ and large $(l>l_{c})$ angular momenta~\cite{Kim1995}.
For small angular momentum $l<l_{c}$ and $\omega=\omega'$, we can approximate $1-\cos(l\theta)\approx(l\theta)^{2}$ such that $[g]=(d+1-\eta)[\theta]=-(\frac{3-\eta}{1+\eta})z$. So, the leading order RG flow equation for $g$ is 
\begin{align}
\frac{dg}{d\ell}={}&-\frac{z(3-\eta)}{1+\eta}g,
\end{align}
which indicates that the interaction part becomes irrelevant, and the system behaves like a Fermi liquid. For large $l>l_{c}$, $\cos(l\theta)$ becomes highly oscillating as a function of $\theta$, so $[g]=(\frac{\eta-1}{\eta+1})z$ and we have,
\begin{align}
\frac{dg}{d\ell}=\frac{z(\eta-1)}{\eta+1}g.
\end{align}
Then the interacting part becomes relevant when $\eta>1$ and we recover non-Fermi liquid behavior. Consequently, even within the quadratic Gaussian action, our effective action can capture the same physics exhibited by the quantum Boltzmann equation for non-Fermi liquids in the collisionless limit, such as the Fermi liquid behavior for $l<l_{c}$ collective modes (between $\Omega\propto q^{(1+\eta)/2}$ and $\Omega\propto q$) and non-Fermi liquid behavior for $l>l_{c}$ collective modes (below $\Omega\propto q^{(1+\eta)/2}$)~\cite{Kim1995}. 

%%%%%%%%%%%%%%%%%%%%%%%%%%%%%%%%%%%%%%%%%%%%%%%%%%%%%%%%%%%%%%%%%%%%%%%%%%%
\section{Notes on the quadratic theory at all loop orders}\label{sec:general}% Notes on general case

In the previous section, we restricted ourselves to the one-loop approximation in the self-energy expansion, $\Sigma^{R}=\Sigma^{R(1)}$. However, it is crucial to verify whether physical results of note, such as the forward scattering cancelation and separation of the angular momentum modes into rough and soft fluctuations, remain valid at higher order in the self-energy expansion or are simply artifacts of the lowest order approximation. To clarify these points, we derive and analyze the Gaussian action in $\phi$ for the general case where $\Sigma^{R}[f]=\sum_{n=1}\Sigma^{R(n)}[f]$. 

Following the same logic as in Sec.~\ref{subsec:quadatic_gaussian}, but now with the full functional expansion of the self-energy, we can write the action in powers of the $\phi$ field.
As shown in Appendix~\ref{app:general_higher}, the $\phi^{0}$ and $\phi^{1}$ order terms again correspond to a constant and a total time derivative, respectively. The quadratic action $S^{(2)}$ can still be split into $S^{(2)}=S^{(2)_{0}}+S^{(2)_{\text{int}}}$, where the non-interacting part $S^{(2)_{0}}$ is the same as in the one-loop approximation (see Eq.~\eqref{eq:cont_quad_bare}). In contrast, the interacting part of the Gaussian action is now 
\begin{align} 
S^{(2)_{\text{int}}}={}&\frac{1}{2}\int_{tx\omega\theta}(\partial_{\omega}f_{0})\dot{\phi}[(\partial_{\omega}\Sigma^{R}_{0})\dot{\phi}-\Delta^{(1)}_{\phi}\Sigma^{R}],\label{eq:general_cont_quad_int}
\end{align}
where $\Sigma^{R}_{0}=\sum_{n=1}\Sigma^{R(n)}[f_0]$. 
In the general case where the full expansion of the self-energy is considered, $\Delta_{\phi}^{(m)}\Sigma^{R(n)}$ can be written in terms of $\bar{\Delta}_{\phi}^{\{m_{i}\}}\Sigma^{R(n)}$, which is defined as
\begin{align}
&\bar{\Delta}_{\phi}^{(m_{1},\cdots,m_{i})}\Sigma^{R(n)}\notag\\&=\int_{\omega_{1}\theta_{1}}\cdots\int_{\omega_{n}\theta_{n}}D_{n}\Delta_{\phi}^{(m_{1})}f_{1}\cdots\Delta_{\phi}^{(m_{i})}f_{i}\prod_{j=i+1}^{n}f_{j}.
\end{align}
$\{m_{j}\}$ is a partition of $m$ which satisfies $\sum_{j=1}^{i}m_{j}=m$. For example, $\Delta_{\phi}^{(2)}\Sigma^{R(n)}=n\bar{\Delta}_{\phi}^{(2)}\Sigma^{R(n)}+\binom{n}{2}\bar{\Delta}_{\phi}^{(1,1)}\Sigma^{R(n)}$.
%%
In Eq.~\eqref{eq:general_cont_quad_int}, $\Delta_{\phi}^{(1)}\Sigma^{R(n)}=n\bar{\Delta}_{\phi}^{(1)}\Sigma^{R(n)}$ such that  $\Delta_{\phi}^{(1)}\Sigma^{R}=\sum_{n=1}n\bar{\Delta}_{\phi}^{(1)}\Sigma^{R(n)}$. Accordingly, $\Sigma^{R}_{0}$ and $\Delta_{\phi}^{(1)}\Sigma^{R}$ now include contributions at all order in the generalized distribution $f$ that are absent from Eq.~\eqref{eq:disc_quad_int}. 


In terms of the Fourier transformed field $\hat{\phi}(\omega,\theta;\Omega,\mathbf{q})$, the interacting part of the quadratic action takes the general form (after discretization of the frequency derivative)
\begin{widetext}
\begin{align}
   S^{(2)_\text{int}}
={}&\frac{1}{2}\int_{\Omega q\omega\theta}F_{0}(\omega,\Omega)(\Sigma^{R}_{0}(\omega+\Omega)-\Sigma^{R}_{0}(\omega))\hat{\phi}(\Omega,\mathbf{q};\omega,\theta)\hat{\phi}(-\Omega,-\mathbf{q};\omega,\theta)\notag\\
&-\sum_{n=1}\frac{n}{2}\int_{\Omega q\omega\theta}\int_{\omega'\theta'}\cdots\int_{\omega_{n-1}\theta_{n-1}}F_{0}(\omega,\Omega)F_{0}(\omega',\Omega)(\prod_{i=1}^{n-1}f_{0}(\omega_{i}))D_{n}(\omega_{1}-\omega,\theta_{1}-\theta;\cdots;\omega'-\omega,\theta'-\theta)\notag\\
&\qquad\qquad\qquad\qquad\qquad\qquad\qquad\times\hat{\phi}(\Omega,\mathbf{q};\omega,\theta)\hat{\phi}(-\Omega,-\mathbf{q};\omega',\theta').\label{eq:general_disc_quad_int}
\end{align}
Using the generic property $D_{n}(\omega_{1},\theta_{1},\cdots;\omega,\theta)=D_{n}(\omega_{1}-\omega,\theta_{1}-\theta;\cdots;\omega_{n}-\omega,\theta_{n}-\theta)$ for the generalized Landau functions,
 $S^{(2)_{\text{int}}}$ can be rewritten as
\begin{align}\label{eq:general_disc_quad_int_symm}
S^{(2)_{\text{int}}}={}&\sum_{n=1}\frac{n}{2}\int_{\Omega q\omega\theta}\int_{\omega'\theta'}\cdots\int_{\omega_{n-1}\theta_{n-1}}F_{0}(\omega,\Omega)F_{0}(\omega',\Omega)(\prod_{i=1}^{n-1}f_{0}(\omega_{i}))D_{n}(\omega_{1}-\omega,\theta_{1}-\theta;\cdots;\omega'-\omega,\theta'-\theta)\notag\\
&\qquad\qquad\qquad\qquad\qquad\qquad\qquad\times\hat{\phi}(\Omega,\mathbf{q};\omega,\theta)[\hat{\phi}(-\Omega,-\mathbf{q};\omega,\theta)-\hat{\phi}(-\Omega,-\mathbf{q};\omega',\theta')].
\end{align}
\end{widetext}
It can then be noted that forward scattering cancellation holds beyond the one-loop approximation since the interacting part of the quadratic action still vanishes near $\omega=\omega'$ and $\theta=\theta'$. Of course, one would expect the forward scattering cancellation to hold at all orders since it can usually be attributed to conservation laws~\cite{kim1994gauge, Kim1995}. Therefore, the explicit demonstration that such a result holds for our EFT at all orders in the self-energy expansion strongly supports its validity.

Similarly, the angular momentum decomposition can be investigated in the general case for the self-energy expansion using the quadratic action \eqref{eq:general_disc_quad_int_symm}. 
Assuming we are working in two dimensions, we can get the angular momentum representation of the general quadratic theory by using Eq.~\eqref{eq:ang_decomp}. 
The non-interacting part of the Gaussian action is once again given by Eq.~\eqref{eq:ang_decomp_quad_bare}, whereas the interacting part is
\begin{widetext}
\begin{align}
S^{(2)_\text{int}}
={}&\sum_{n=1}\frac{n}{2}\sum_{l}\int_{\Omega q}\int_{\omega\omega'}\int_{\delta\theta'}\cdots\int_{\omega_{n-1}\delta\theta_{n-1}}F_{0}(\omega,\Omega)F_{0}(\omega',\Omega)(\prod_{i=1}^{n-1}f_{0}(\omega_{i}))D_{n}(\omega_{1}-\omega,\delta\theta_{1};\cdots;\omega'-\omega,\delta\theta')\notag\\
&\qquad\qquad\qquad\qquad\times\tilde{\phi}(-\Omega,-\mathbf{q};\omega,l)[\tilde{\phi}(\Omega,\mathbf{q};\omega,l)-\cos(l\theta)\tilde{\phi}(\Omega,\mathbf{q};\omega',l)],
%%%%
\end{align}
\end{widetext}
where $\delta\theta'=|\theta'-\theta|$ and $\delta\theta_{i}=|\theta_{i}-\theta|$. For small angular momentum $l$, $S^{(2)_\text{int}}$ vanishes, and the system thus exhibits Fermi liquid behavior. On the other hand, the cancellation does not happen for the large angular momentum, and the interacting part of the quadratic action thus leads to the non-Fermi liquid-like rough fluctuations. This is analogous to the case discussed in Sec.~\ref{subsec:quadratic_angular} within the one-loop approximation. Therefore, the conclusions regarding the angular momentum decomposition of the EFT obtained in the one-loop approximation also hold in the general case.

%%%%%%%%%%%%%%%%%%%%%%%%%%%%%%%%%%%%
\section{Discussion}\label{sec:discussion}
This work proposed a bosonized effective field theory for non-Fermi liquid. In our scheme, the fluctuating bosonic field $\phi$ parametrizes the Fermi surface deformation and may be interpreted as frequency-dependent local chemical potential variations in momentum space. The proposed action reproduces the quantum Boltzmann equation in the collisionless limit as its equation of motion. We obtained the Gaussian theory by expanding the non-linear action at quadratic order for $\phi$. The Gaussian action shows the usual Fermi liquid behavior for forward scattering due to a cancellation between the self-energy and vertex correction as seen from the Ward identity~\cite{kim1994gauge, Kim1995}. We computed the Gaussian contribution to the two-point density correlation function perturbatively and described how they relate to the loop corrections in the equilibrium theory of a Fermi surface interacting with gapless bosons through a Yukawa coupling. In the angular momentum decomposition of the Gaussian action, we demonstrated that it exhibits Fermi liquid-like behavior for small angular momentum and non-Fermi liquid behavior for large angular momentum. These results obtained only at the quadratic level are consistent with the previous studies of the quantum Boltzmann equation for non-Fermi liquid in the collisionless limit~\cite{Kim1995, Mandal2022}.
We have further demonstrated that these results are valid whether the self-energy is expanded at the one-loop order or considered in its most general functional expansion form.

A key point of our formalism is the use of the generalized distribution $f(\omega,\theta;t,\mathbf{x})=-i \int \frac{d \xi_\mathbf{p}}{2 \pi} G^{<}(\mathbf{p}, \omega ; \mathbf{x}, t)$ rather than the semiclassical phase space distribution used to describe the dynamics of ordinary Fermi liquids. Because of its dependence on the relative frequency $\omega$, the fluctuations of $f$ need to be parametrized by the generalized canonical transformation~\eqref{eq:canonical_transformatioon_quasiclassical_distribution} where the generalized Poisson bracket, with its additional time and frequency derivatives, is used. This novel approach allows us to effectively capture the intricate frequency dependence of the self-energy, which serves as a fundamental characteristic of non-Fermi liquid behavior. This explicit consideration of the frequency through the use of the generalized distributions and consideration of the self-energy is a crucial difference from the previous literature.

In this article, we have primarily focused on the quadratic action. Even at the Gaussian level, we have shown that our EFT reproduces most of the physics of the quantum Boltzmann equation, such as Fermi liquid and non-Fermi liquid behaviors in certain limits, which are consistent with previous results in the literature~\cite{Kim1995,Mandal2022}. Although much of the known physics can be reproduced at the quadratic order in $\phi$, higher-order effects in $\phi$ might still play an essential role in understanding physical observations. Our formalism provides an explicit way to consider these higher-order corrections systematically (see Appendices~\ref{app:higher} and~\ref{app:general_higher}). Non-linear effects and non-linear responses such as $n$-point density correlation function~\cite{Son2022} should thus be carefully considered in future studies.


When studying the proposed EFT, it was noted that the self-energy contribution to the action can be rewritten in terms of a functional $\Phi[f]$ (see Eq.~\eqref{eq:definition_generalized_luttinger-Ward}) that satisfies $\delta\Phi[f]/\delta f=\Sigma^{R}[f]$ --- in analogy with the conventional Luttinger-Ward functional. This intriguing observation requires clarification on whether the correspondence is fundamental or a mere coincidence. We expect further investigations of the functional to unveil a more concrete relationship between $\Phi[f]$ and the Luttinger-Ward functional or potentially other generating functionals of the equilibrium theory studied previously.

Our work further opens the doors to exciting and important generalizations. For instance, we have only considered the spinless Fermi surface. A natural extension will be the bosonization of the spinful (multiflavor) non-Fermi liquid. The quadratic action will be similar to the spinless case. Still, higher-order contributions to the action for $\phi$, such as the cubic terms, can differ significantly due to the nontrivial Lie algebra associated with the internal symmetry of the fermions~\cite{Son2022}. Therefore, analysis of higher-order terms in the action will become particularly important in studying spinful (multiflavor) Fermi surfaces.


\begin{acknowledgements}
We thank Daniel J. Schultz, Omid Tavakol, Hart Goldman, T.~Senthil, Yi-Hsien Du and D.~T.~Son for insightful discussions.  We acknowledge support from the Natural Sciences and Engineering Research Council of Canada (NSERC) and the Centre of Quantum Materials at the University of Toronto. F.D. is further supported by the Vanier Canada Graduate Scholarship. YBK is also supported by the Guggenheim Fellowship from the John Simon Guggenheim Memorial Foundation and the Simons Fellowship from the Simons Foundation. This work was partly performed at the Aspen Center for Physics, which is supported by National Science Foundation grant PHY-2210452. YBK acknowledges the support of the Advanced Study Group on ``Entanglement and Dynamics of Quantum Matter" at the Center for Theoretical Physics of Complex Systems in the Institute for Basic Science, where a part of the current work was done.
\end{acknowledgements}

%%%%%%%%%%%%%%%%%%%%%%%%%%%%%%%%%%
\appendix
%%%%%%%%%%%%%%%%%%%%%%%%%%%%%%%%%%


%%%%%%%%%%%%%%%%%%%%%%%%%%%%%%%%%%
\section{Non-equilibrium field theory} \label{app:Non-equilibirum_field_theory}

To derive transport equations of metals, one has to consider a non-equilibrium situation. We here briefly review non-equilibrium field theory. The usual prescription is to work a closed-time Keldysh contour $\mathcal{C}$ (i.e., $\mathcal{C}$ describes a path from the distant past to the time of interest and back). The non-equilibrium Dyson's equation can be written as~\cite{kita2010introduction, rammer1986quantum, rammer2011quantum}
\begin{align} \label{eq:Dyson_equation_matrix_GF}
    \left(i  \frac{\partial}{\partial t_1} - H_0\right) \check{G}(1,2)-\int d3 \check{\Sigma}(1,3) \check{G}(3,2) = \mathds{1}\delta(1,2)
\end{align}
where $1\equiv (t_1, \mathbf{x}_1)$, $\delta(1,2)\equiv \delta(t_1-t_2)\delta^{(d)}(\mathbf{x}_1-\mathbf{x}_2)$,  $H_0$ is the non-interacting Hamiltonian, and the matrix Green's function and self-energy are
\begin{align}  \label{eq:Definition_matrix_GF_and_SE}
\check{G}=\left[\begin{array}{cc}G^{\mathcal{T}} & -G^{<} \\ G^{>} & -G^{\overline{\mathcal{T}}} \end{array}\right] \text{ and } \check{\Sigma}=\left[\begin{array}{ll} \Sigma^{\mathcal{T}} & -\Sigma^{<} \\ \Sigma^{>} & -\Sigma^{\overline{\mathcal{T}}} \end{array}\right]
\end{align}
with
\begin{subequations} \label{eq:Definition_of_all_GFs}
\begin{align} 
G^{>}\left(1, 2\right)= & -i\left\langle\psi\left(1\right) \psi^{\dagger}\left(2\right)\right\rangle \\ 
G^{<}\left(1, 2\right)= & i\left\langle\psi^{\dagger}\left(2\right) \psi\left(1\right)\right\rangle \\ 
G^{\mathcal{T}}\left(1, 2\right)= & \Theta\left(t_1-t_2\right) G^{>}\left(1, 2\right) \notag\\
&\hspace{1cm}+\Theta\left(t_2-t_1\right) G^{<}\left(1, 2\right) \\ 
G^{\overline{\mathcal{T}}}\left(1, 2\right)= & \Theta\left(t_2-t_1\right) G^{>}\left(1, 2\right) \notag\\
&\hspace{1cm}
 + \Theta\left(t_1-t_2\right) G^{<}\left(1, 2\right).
\end{align}
\end{subequations}

We note once again that Green's functions will not have spatial or temporal translation symmetry in the general non-equilibrium case. This implies that they will depend on both absolute (i.e., $\mathbf{x}=(\mathbf{x}_1+\mathbf{x}_2)/2$ and $t=(t_1+t_2)/2$) and relative coordinates (i.e., $\mathbf{x}_{\mathrm{rel}}=\mathbf{x}_{2} - \mathbf{x}_1$ and $t_{\mathrm{rel}}=t_{2}-t_{1}$). To derive equations of motions, it is convenient to work in the Wigner representation
\begin{align}
    \check{G}\left(\omega, \mathbf{p};  t, \mathbf{x} \right)  &= \mathcal{W}\left[\check{G}\left(t_{\text{rel}}, \mathbf{x}_{\text{rel}};  t, \mathbf{x} \right)\right]  \notag \\
    &= \int d^{d} x_{\text{rel}} dt_{\text{rel}} \check{G}(1,2) e^{-i\left(\mathbf{p} \cdot \mathbf{x}_{\text{rel}}- \omega t_{\text{rel}}\right)},
\end{align}
which is a Fourier transform with respect to the relative coordinates. To obtain the reformulation of Dyson's equation in this representation, one needs to apply the Wigner transformation to Eq.~\eqref{eq:Dyson_equation_matrix_GF} which contains matrix products. The Wigner transformation of a product of function is formally defined as the Groenewold-Moyal product~\cite{moyal1949quantum, groenewold1946principles} but can be approximated at first order in the gradient expansion by
\begin{align}
    \mathcal{W}\left[B C\right] \approx B C + \frac{i}{2} \left[B, C\right],
\end{align}
where we have introduced the generalized Poisson bracket defined in Eq.~\eqref{eq:generalized_Poisson_bracket}. Therefore, in the first-order gradient approximation of the Wigner representation, Dyson's equation leads to
\begin{equation} \label{eq:Dyson_equation_at_first_order_in_Wigner_representation}
    [G_{0}^{-1}-\Re\Sigma^{R}, G^{<}] - [\Sigma^{<}, \Re G^{R}] = G^{>} \Sigma^{>} - G^{<} \Sigma^{<},
\end{equation}
where we introduced $(G_0)^{-1} = \omega - \xi_{\mathbf{p}}$  with $\xi_{\mathbf{p}}=\epsilon(\mathbf{p})-\mu$ denoting the spectrum of $H_0$ (assumed to be independent of absolute position $\mathbf{x}$ and time $t$) and the retarded Green's function 
\begin{equation}
    G^{R}(1,2) = \theta(t_1 - t_2) \left( G^{>}(1,2) - G^{<}(1,2) \right)
\end{equation}
and its self-energy. The right-hand side of Eq.~\eqref{eq:Dyson_equation_at_first_order_in_Wigner_representation} is commonly referred to as the collision integral. 

%%%%%%%%%%%%%%%%%%%%%%%%%%%%%%%%%%
\section{Useful relations}\label{app:useful}
In this section, we present derivations of some useful relations used in the main text.

\noindent \underline{1) Vanishing of $\int \frac{d\xi_{\mathbf{p}}}{2\pi}\text{Re}G^{R}$}
\begin{align}
    \int&\frac{d\xi_{\mathbf{p}}}{2\pi}\text{Re}G^{R}(\omega',\mathbf{p};t,\mathbf{x})\notag\\
    ={}&\int\frac{d\omega'}{2\pi}\mathcal{P}\frac{[1-f(\omega',\theta;t,\mathbf{x})]+f(\omega',\theta;t,\mathbf{x})}{\omega-\omega'}\notag\\
    ={}&\int\frac{d\omega'}{2\pi}\mathcal{P}\frac{1}{\omega-\omega'}=0,
\end{align}
where $\mathcal{P}$ means the principal value, and
\begin{align}
    \text{Re}&G^{R}(\omega',\mathbf{p};t,\mathbf{x})\notag\\
    ={}&-\int\frac{d\omega'}{2\pi}\mathcal{P}\frac{G^{>}(\omega',\mathbf{p};t,\mathbf{x})-G^{<}(\omega',\mathbf{p};t,\mathbf{x})}{\omega-\omega'},\\
    \int&\frac{d\xi_{\mathbf{p}}}{2\pi}[-iG^{<}(\omega,\mathbf{p};t,\mathbf{x})]={}f(\omega,\theta;t,\mathbf{x}),\\
    \int&\frac{d\xi_{\mathbf{p}}}{2\pi}[iG^{>}(\omega,\mathbf{p};t,\mathbf{x})]={}1-f(\omega,\theta;t,\mathbf{x}).
\end{align}

\noindent \underline{2) $\int_{t\omega}\braket{\Delta_{\phi}^{(m)}f,\Delta_{\phi}^{(n)}\Sigma^{R}}=\int_{t\omega}\braket{\Delta_{\phi}^{(n)}f,\Delta_{\phi}^{(m)}\Sigma^{R}}$}
\begin{align}
\int_{t\omega}&\braket{\Delta_{\phi}^{(m)}f,\Delta_{\phi}^{(n)}\Sigma^{R}}\notag\\
={}&\int_{tx\omega p}(\Delta_{\phi}^{(m)}f)(\Delta_{\phi}^{(n)}\Sigma^{R})\notag\\
={}&\int_{tx\omega p}\int_{\omega'\theta'}(\Delta_{\phi}^{(m)}f(\omega))D_{1}(\omega'-\omega,\theta'-\theta)(\Delta_{\phi}^{(n)}f(\omega'))\notag\\
={}&\int_{tx\omega p}[\int_{\omega'\theta'}(\Delta_{\phi}^{(m)}f(\omega'))D_{1}(\omega'-\omega,\theta'-\theta)](\Delta_{\phi}^{(n)}f(\omega))\notag\\
={}&\int_{tx\omega p}(\Delta_{\phi}^{(n)}f)(\Delta_{\phi}^{(m)}\Sigma^{R})\notag\\
={}&\int_{t\omega}\braket{\Delta_{\phi}^{(n)}f,\Delta_{\phi}^{(m)}\Sigma^{R}},
\end{align}
where $\Delta_{\phi}^{(n)}\Sigma^{R}=\bar{\Delta}_{\phi}^{(n)}\Sigma^{R(1)}$ in the one-loop approximation (i.e., $\Sigma[f]=\Sigma^{(1)}[f]$) and we used the fact that $D_{1}$ does not change its sign under the change of variables~\eqref{eq:Change_of_variables_Landau_parameters}. Note that this identity is only valid in the one-loop approximation.

\begin{widetext}
\section{Equation of motion}\label{app:general_action}
Let us consider the following self-energy contribution which contains $n$ factors of $f$
\begin{align}
\Sigma^{R(n)}={}&\int_{1}\int_{2}\ldots\int_{n}D_{n}(\omega_{1},\omega_{2},\ldots,\omega_{n})f_{1}f_{2}\ldots f_{n},
\end{align}
where $f_{i}=f(\omega_{i},\theta_{i})=U_{i}f_{0}(\omega_{i})U_{i}^{-1}$, $U_{i}=e^{-\phi(\omega_{i},\theta_{i})}$, $\int_{i}=\int_{\omega_{i}\theta_{i}}$, and we assume that $D_{n}(\omega_{1},\theta_{1},\ldots;\omega_{n},\theta_{n})$ is invariant under the change of variables~\eqref{eq:Change_of_variables_Landau_parameters}. Note that the $t$ and $\mathbf{x}$ dependence of $\Sigma^{R}$, $f$, and $\phi$ are not indicated for simplicity. Then, the self-energy changes under $U\rightarrow e^{-\alpha}U$, where $\alpha$ is an infinitesimal variation, according to
\begin{align}
 U^{-1}{}&\Sigma^{R(n)} U= U^{-1}\int_{1}\int_{2}\ldots\int_{n}( D_{n}f_{1}f_{2}\ldots f_{n})U\notag\\
={}& U^{-1}\int_{1}\int_{2}\ldots\int_{n}( D_{n}(U_{1}f_{0}U_{1}^{-1})(U_{2}f_{0}U_{2}^{-1})\ldots(U_{n}f_{0}U_{n}^{-1}))U\notag\\
\rightarrow& U^{-1}e^{\alpha}\int_{1}\int_{2}\ldots\int_{n}( D_{n}(e^{-\alpha_{1}}U_{1}f_{0}U_{1}^{-1}e^{\alpha_{1}})(e^{-\alpha_{2}}U_{2}f_{0}U_{2}^{-2}e^{\alpha_{2}})\ldots(e^{-\alpha_{n}}U_{n}f_{0}U_{n}^{-1}e^{\alpha_{n}}))e^{-\alpha}U\notag\\
=& U^{-1}\Bigl(\int_{1}\int_{2}\ldots\int_{n}D_{n}(e^{-\alpha_{1}}f_{1}e^{\alpha_{1}})(e^{-\alpha_{2}}f_{2}e^{\alpha_{2}})\ldots(e^{-\alpha_{n}}f_{n}e^{\alpha_{n}})\notag\\
&\qquad\qquad\qquad\qquad\qquad+[\alpha,\int_{1}\int_{2}\ldots\int_{n}D_{n}(e^{-\alpha_{1}}f_{1}e^{\alpha_{1}})(e^{-\alpha_{2}}f_{2}e^{\alpha_{2}})\ldots(e^{-\alpha_{n}}f_{n}e^{\alpha_{n}})]+\ldots\Bigr)U\notag\\
=& U^{-1}\Bigl(\int_{1}\int_{2}\ldots\int_{n}D_{n}(\prod_{i=1}^{n}f_{i}-n(\prod_{i=1}^{n-1}f_{i})[\alpha_{n},f_{n}]_{n}+\ldots)+[\alpha,\Sigma^{R(n)}+\ldots]+\ldots\Bigr)U,
\end{align}
where $[A,B]_{n}$ is the generalized Poisson bracket including the derivatives with respect to $\omega_{n}$ and $\theta_{n}$. It can then be noticed that for the contribution to the self-energy which contains $n$ generalized distribution function there is $n$ additional contributions to the linear order variation. 

Considering the general form of the action,
\begin{align*}
S=\int_{t\omega}\braket{f,\omega-\epsilon-\sum_{n=1}\tfrac{1}{n+1}\Sigma^{R(n)}}.
\end{align*}
we have
\begin{align}
\delta S={}&\int_{t\omega}\braket{f_{0},-U^{-1}[\omega-\xi_{\mathbf{p}},\alpha]U-\frac{1}{n+1}U^{-1}\Bigl(n\int_{1}\int_{2}\ldots\int_{n}D_{n}(\prod_{i=1}^{n-1}f_{i})[\alpha_{n},f_{n}]_{n}-[\Sigma^{R(n)},\alpha]\Bigl)U}\notag\\
={}&\int_{t\omega}\braket{[\omega-\xi_{\mathbf{p}}-\tfrac{1}{n+1}\Sigma^{R(n)},f],\alpha}-\sum_{n=1}\frac{n}{n+1}\int_{t\omega}\braket{\int_{1}\int_{2}\ldots\int_{n}[D_{n}f(\prod_{i=1}^{n-1}f_{i}),f_{n}]_{n},\alpha_{n}}\notag\\
={}&\int_{t\omega}\braket{[\omega-\xi_{\mathbf{p}}-\sum_{n=1}\tfrac{1}{n+1}\Sigma^{R(n)},f],\alpha}-\sum_{n=1}\frac{n}{n+1}\int_{t\omega}\braket{\int_{1}\int_{2}\ldots\int_{n}[D_{n}(\prod_{i=1}^{n}f_{i}),f],\alpha}\notag\\
={}&\int_{t\omega}\braket{[\omega-\xi_{\mathbf{p}}-\sum_{n=1}\tfrac{1}{n+1}\Sigma^{R(n)},f],\alpha}-\sum_{n=1}\frac{n}{n+1}\int_{t\omega}\braket{[\Sigma^{R(n)},f],\alpha}\notag\\
={}&\int_{t\omega}\braket{[\omega-\xi_{\mathbf{p}}-\Sigma^{R},f],\alpha},
\end{align}
so that the equation of motion is
\begin{align}
[\omega-\epsilon-\Sigma^{R},f]=0
\end{align}
as mentioned in the main text. If $D_{n}$ is not invariant under the change of variables~\eqref{eq:Change_of_variables_Landau_parameters}, the coefficient of $\Sigma^{R(n)}$ in the effective action can differ.



\section{The generalized Luttinger-Ward functional \texorpdfstring{$\Phi[f]$}{}}\label{app:LW}
The self-energy contribution to our effective action including $n$ factors of the generalized distribution function under some assumptions (see Appendix~\ref{app:general_action}) can be written as
\begin{align}
\int_{\omega\theta}{}&f(\omega,\theta)\sum_{n=1}\frac{1}{n+1}\Sigma^{R(n)}(\omega,\theta)\notag\\
=&\int_{\omega\theta}f(\omega,\theta)\left( \frac{1}{2!}\int_{\omega_{1}\theta_{1}}\frac{\delta\Sigma^{R}}{\delta f(\omega_{1},\theta_{1})}f(\omega_{1},\theta_{1})
+\frac{1}{3!}\int_{\omega_{1}\theta_{1}}\int_{\omega_{2}\theta_{2}}\frac{\delta^{2}\Sigma^{R}}{\delta f(\omega_{1},\theta_{1})\delta f(\omega_{2},\theta_{2})}f(\omega_{1},\theta_{1})f(\omega_{2},\theta_{2})+\ldots\right)\notag\\
={}&\Phi[f].
\end{align}
Taking the derivative of the functional $\Phi[f]$ with respect to $f$, we have
\begin{align}
    \frac{\delta\Phi[f]}{\delta f(\omega,\theta)}
    ={}&\int_{\omega_{1}\theta_{1}}\frac{\delta\Sigma^{R}}{\delta f(\omega_{1},\theta_{1})}f(\omega_{1},\theta_{1})+\frac{1}{2!}\int_{\omega_{1}\theta_{1}}\int_{\omega_{1}\theta_{1}}\frac{\delta^{2}\Sigma^{R}}{\delta f(\omega_{1},\theta_{1})\delta f(\omega_{2},\theta_{2})}f(\omega_{1},\theta_{1})f(\omega_{2},\theta_{2})+\ldots={}\Sigma^{R}[f].
\end{align}
So the functional has the property $\frac{\delta\Phi[f]}{\delta f}=\Sigma^{R}[f]$, which is analogous to the Luttinger-Ward functional. Then, following this analogy, we may consider the prefactor $1/(n+1)$ in front of the self-energy contribution in the effective action Eq.~\eqref{eq:action} as the inverse of the symmetry factor for the corresponding one-particle irreducible diagram including the $n$ generalized distribution functions. 
Thus, by using this generalized Luttinger-Ward functional, the effective action can be written as
\begin{align}
    S={}&\int_{t\omega}\braket{f,\omega-\epsilon}-\int_{tx}\Phi[f].
\end{align}


\section{Comments on non-linear corrections in the one-loop approximation}\label{app:higher}
We can also obtain higher-order terms in $\phi$ beyond the quadratic action, which gives us non-linear corrections.
From Eq.~\eqref{eq:effective_expand}, the general form of $n$-th order action is given by
\begin{align}
S^{(n)}={}&\int_{t\omega}\braket{\Delta_{\phi}^{(n)}f,\omega-\epsilon(p)-\tfrac{1}{2}\Sigma^{R}_{0}(\omega)}-\int _{t\omega}\sum_{i=1}^{n}\braket{\Delta_{\phi}^{(n-i)}f,\tfrac{1}{2}\Delta_{\phi}^{(i)}\Sigma^{R}}.
\end{align}

For example, the cubic action is
\begin{align}
    S^{(3)}={}&\int dtd\omega\braket{\Delta_{\phi}^{(3)}f,\omega-\epsilon(p)-\Sigma^{R}_{0}(\omega)}-\int dtd\omega\braket{\Delta_{\phi}^{(2)}f,\Delta_{\phi}^{(1)}\Sigma^{R}}.
\end{align}
The non-interacting part of this cubic term can be expressed in terms of $\phi$ as
\begin{align}
S^{(3)_{0}}={}&\int_{t\omega}\braket{\Delta^{(3)}_{\phi}f,\omega-\xi_{p}}\notag\\
={}&\frac{1}{6}\int_{tx\omega\theta}(\partial_{\omega}^{2}f_{0})\phi\dot{\phi}\ddot{\phi}
+\frac{1}{6}\int_{tx\omega\theta}(\partial_{\omega}f_{0})\dot{\phi}[p_{F}^{-1}\epsilon'(\mathbf{s}_{\theta}\cdot\nabla_{\mathbf{x}}\phi)^{2}+\epsilon''(\mathbf{n}_{\theta}\cdot\nabla_{\mathbf{x}}\phi)^{2}]\notag\\
&-\frac{1}{4}\int_{tx\omega\theta}\dot{\phi}^{2}(\partial_{\omega}f_{0})(\nabla_{\mathbf{p}}\epsilon)\cdot(\nabla_{\mathbf{x}}\partial_{\omega}\phi)\notag
\\
&-\frac{1}{6}\int_{tx\omega\theta}[(\nabla_{\mathbf{x}}\phi)\cdot(\nabla_{\mathbf{p}}\dot{\phi})-(\nabla_{\mathbf{p}}\phi)\cdot(\nabla_{\mathbf{x}}\dot{\phi})](\partial_{\omega}f_{0})(\nabla_{\mathbf{p}}\epsilon)\cdot(\nabla_{\mathbf{x}}\phi),\label{eq:cont_cubic_bare}
\end{align}
and the interacting part is
\begin{align}
S^{(3)_{\text{int}}}={}&-\int_{t\omega}\braket{\Delta^{(3)}_{\phi}f,\Sigma^{R}_{0}}-\int_{t\omega}\braket{\Delta^{(2)}_{\phi}f,\Delta^{(1)}_{\phi}\Sigma^{R}}\notag\\
={}&-\frac{1}{6}\int_{tx\omega\theta}[(\partial_{\omega}^{2}f_{0})(\partial_{\omega}\Sigma^{R}_{0})-(\partial_{\omega}f_{0})(\partial_{\omega}^{2}\Sigma^{R}_{0})]\phi\dot{\phi}\ddot{\phi}\notag\\
&+\frac{1}{2}\int_{tx\omega\theta}(\partial_{\omega}f_{0})\dot{\phi}[\dot{\phi}(\partial_{\omega}\Delta^{(1)}_{\phi}\Sigma^{R})-(\partial_{\omega}\phi)(\partial_{t}\Delta^{(1)}_{\phi}\Sigma^{R})]\notag\\
&+\frac{1}{2}\int_{tx\omega\theta}[(\nabla_{\mathbf{x}}\phi)\cdot(\nabla_{\mathbf{p}}\dot{\phi})-(\nabla_{\mathbf{p}}\phi)\cdot(\nabla_{\mathbf{x}}\dot{\phi})](\partial_{\omega}f_{0})\Delta^{(1)}_{\phi}\Sigma^{R}.\label{eq:cont_cubic_int}
\end{align}

The discretized non-interacting cubic action is
\begin{align}
S^{(3)_{0}}
={}&\frac{i}{6}\int_{\Omega q}'\int_{\omega\theta}[F_{0}(\omega,\Omega_{1})+F_{0}(\omega,-\Omega_{1})]\Omega_{2}\hat{\phi}_{1}\hat{\phi}_{2}\hat{\phi}_{3}\notag\\
&+\frac{i}{6}\int_{\Omega q}'\int_{\omega\theta}F_{0}(\omega,\Omega_{1})[p_{F}^{-1}\epsilon'(\mathbf{s}_{\theta}\cdot\mathbf{q}_{2})(\mathbf{s}_{\theta}\cdot\mathbf{q}_{3})+\epsilon''(\mathbf{n}_{\theta}\cdot\mathbf{q}_{2})(\mathbf{n}_{\theta}\cdot\mathbf{q}_{3})]\hat{\phi}_{1}\hat{\phi}_{2}\hat{\phi}_{3}\notag\\
&+\frac{i}{4}\int_{\Omega q}'\int_{\omega\theta}F_{0}(\omega,\Omega_{1})\hat{\phi}_{1}(v_{F}\mathbf{n}_{\theta}\cdot\mathbf{q}_{2})[\hat{\phi}_{2}(\omega+\Omega_{3})-\hat{\phi}_{2}(\omega)]\hat{\phi}_{3}\notag
\\
&-\frac{i}{6}\int_{\Omega q}'\int_{\omega\theta}F_{0}(\omega,\Omega_{1})[p_{F}^{-1}(\mathbf{q}_{2}\cdot\nabla_{\theta}\hat{\phi}_{1})\hat{\phi}_{2}-p_{F}^{-1}\hat{\phi}_{1}(\mathbf{q}_{1}\cdot\nabla_{\mathbf{\theta}}\hat{\phi}_{2})](v_{F}\mathbf{n}_{\theta}\cdot\mathbf{q}_{3})\hat{\phi}_{3},\label{eq:disc_cubic_bare}
\end{align}
and the discretized interacting cubic action is
\begin{align}
S^{(3)_{\text{int}}}={}&-\frac{i}{6}\int_{\Omega q}'\int_{\omega\theta}\int_{\omega'\theta'}[F_{0}(\omega',\Omega_{1})(F_{0}(\omega,\Omega_{2})+F_{0}(\omega,-\Omega_{2})-F_{0}(\omega,\Omega_{1})(F_{0}(\omega',\Omega_{2})+F_{0}(\omega',-\Omega_{2})]D_{1}(\omega'-\omega,\theta'-\theta)\hat{\phi}_{1}\hat{\phi}_{2}\hat{\phi}_{3}\notag\\
&+\frac{i}{2}\int_{\Omega q}'\int_{\omega\theta}\int_{\omega'\theta'}F_{0}(\omega,\Omega_{1})F_{0}(\omega',\Omega_{2})(D_{1}(\omega'-\omega-\Omega_{3},\theta'-\theta)-D_{1}(\omega'-\omega,\theta'-\theta))\hat{\phi}_{1}\hat{\phi}_{2}\hat{\phi}_{3}'\notag\\
&-\frac{i}{2}\int_{\Omega q}'\int_{\omega\theta}\int_{\omega'\theta'}F_{0}(\omega,\Omega_{1})F_{0}(\omega',\Omega_{3})D_{1}(\omega'-\omega,\theta'-\theta)\hat{\phi}_{1}(\hat{\phi}_{2}(\omega+\Omega_{3})-\hat{\phi}_{2}(\omega))\hat{\phi}_{3}'\notag\\
&-\frac{i}{2}\int_{\Omega q}'\int_{\omega\theta}\int_{\omega'\theta'}F_{0}(\omega,\Omega_{1})F_{0}(\omega',\Omega_{3})[p_{F}^{-1}(\mathbf{q}_{2}\cdot\nabla_{\theta}\hat{\phi}_{1})\hat{\phi}_{2}-p_{F}^{-1}\hat{\phi}_{1}(\mathbf{q}_{1}\cdot\nabla_{\theta}\hat{\phi}_{2})]D_{1}(\omega'-\omega,\theta'-\theta)\hat{\phi}_{3}'
\end{align}
where $\int_{\Omega q}'=\int\frac{d\Omega_{1} d^{d}q_{1}}{(2\pi)^{d+1}}\frac{d\Omega_{2} d^{d}q_{2}}{(2\pi)^{d+1}}\frac{d\Omega_{3} d^{d}q_{3}}{(2\pi)^{d+1}}\delta(\Omega_{1}+\Omega_{2}+\Omega_{3})\delta(\mathbf{q}_{1}+\mathbf{q}_{2}+\mathbf{q}_{3})$, and we used that $(\partial_{\omega}^{2}f_{0})\Omega^{2}\approx[f_{0}(\omega+\Omega)-f_{0}(\omega)]+[f_{0}(\omega-\Omega)-f_{0}(\omega)]$, $(\partial_{\omega}f_{0})\Omega\approx f_{0}(\omega+\Omega)-f_{0}(\omega)$, and $(\partial_{\omega}\hat{\phi}(\omega))\Omega\approx \hat{\phi}(\omega+\Omega)-\hat{\phi}(\omega)$. Note that we have also used the shorthand notation $F_{0}(\omega,\Omega)=f_{0}(\omega+\Omega)-f_{0}(\omega)$, $\hat{\phi}_{n}\equiv\hat{\phi}(\omega,\theta;\Omega_{n},\mathbf{q}_{n})$,  $\hat{\phi}'_{n}\equiv\hat{\phi}(\omega',\theta';\Omega_{n},\mathbf{q}_{n})$, and $\hat{\phi}(\omega)=\hat{\phi}(\omega,\theta;\Omega,\mathbf{q})$.
\end{widetext}


%%%%%%%%%%%%%%%%%%%%%%%%%%%%%%
\section{Quadratic action for \texorpdfstring{$\phi$ in the one-loop approximation}{}}\label{app:quadratic}
The Gaussian action in the one-loop approximation is given by
\begin{align}
    S^{(2)}={}&\int_{t\omega}\braket{\Delta_{\phi}^{(2)}f,\omega-\xi_{\mathbf{p}}-\Sigma^{R}_{0}}%\notag\\&
    -\int_{t\omega}\braket{\Delta_{\phi}^{(1)}f,\tfrac{1}{2}\Delta_{\phi}^{(1)}\Sigma^{R}}.
\end{align}
It can be computed as follows. The first part is
\begin{align}
    \int_{t\omega}&\braket{\Delta_{\phi}^{(2)}f,\omega-\xi_{\mathbf{p}}-\Sigma^{R}_{0}}\notag\\
    ={}&\frac{1}{2}\int_{t\omega}\braket{[\phi,[\phi,f_{0}]],\omega-\xi_{\mathbf{p}}-\Sigma^{R}_{0}}\notag\\
    ={}&-\frac{1}{2}\int_{t\omega}\braket{[\phi,f_{0}],[\phi,\omega-\xi_{\mathbf{p}}-\Sigma^{R}_{0}]}\notag\\
    ={}&\frac{1}{2}\int_{t\omega}\braket{(\partial_{\omega}f_{0})\dot{\phi},-\dot{\phi}-(\nabla_{\mathbf{p}}\xi_{\mathbf{p}})\cdot(\nabla_{\mathbf{x}}\phi)+(\partial_{\omega}\Sigma^{R}_{0})\dot{\phi}]}\notag\\
    ={}&-\frac{1}{2}\int_{t\omega xp}(\partial_{\omega}f_{0})\dot{\phi}[\dot{\phi}+(\nabla_{\mathbf{p}}\xi_{\mathbf{p}})\cdot(\nabla_{\mathbf{x}}\phi)]\notag\\
    &+\frac{1}{2}\int_{t\omega xp}(\partial_{\omega}f_{0})\dot{\phi}^{2}(\partial_{\omega}\Sigma^{R}_{0}),
\end{align}
and second part is
\begin{align}
    %
    \int_{t\omega}&\braket{\Delta_{\phi}^{(1)}f,\tfrac{1}{2}\Delta_{\phi}^{(1)}\Sigma^{R}}\notag\\
    ={}&-\frac{1}{2}\int_{t\omega}\braket{[\phi,f_{0}],\Delta_{\phi}^{(1)}\Sigma^{R}}\notag\\
    ={}&\frac{1}{2}\int_{t\omega}\braket{(\partial_{\omega}f_{0})\dot{\phi},\Delta_{\phi}^{(1)}\Sigma^{R}}
    ={}\frac{1}{2}\int_{tx\omega\theta}(\partial_{\omega}f_{0})\dot{\phi}\Delta_{\phi}^{(1)}\Sigma^{R},
\end{align}
where we used the following relation,
\begin{align}
    \int_{t\omega}&\braket{A,B}\notag\\={}&\mathcal{C}\int\frac{dtd\omega d^{d}xd^{d}p}{(2\pi)^{d+1}}AB\\
    ={}&\frac{N(\mu)}{\left(\int\frac{d\epsilon}{2\pi}N(\epsilon)\right)}\int\frac{dtd\omega d^{d}xd^{d-1}\theta}{(2\pi)^{d}}\int\frac{d\epsilon}{2\pi}N(\epsilon)AB\\
    ={}&N(\mu)\int\frac{dtd\omega d^{d}xd^{d-1}\theta}{(2\pi)^{d}}AB=\int_{tx\omega\theta}AB,
\end{align}
under the assumption that $AB$ does not depend on $\epsilon$, and the approximation $\xi_{\mathbf{p}}\approx v_{F}(\theta) p_{\parallel}$, so that $\nabla_{\mathbf{p}}\xi_{\mathbf{p}}$ does not depend on $\xi_{\mathbf{p}}$ (e.g, $\nabla_{\mathbf{p}}\xi_{\mathbf{p}}=v_{F}\mathbf{n}_{\theta}$ in the isotropic limit). 
Note that $\Sigma^{R}_{0}$ only depends on $\omega$.
Then, as a result, we can obtain Eq.~\eqref{eq:cont_quad_bare} and \eqref{eq:cont_quad_int} in the main text.
Here, if we perform the Fourier transformation for $\phi$, we can easily obtain Eq.~\eqref{eq:disc_quad_bare} in the main text. In the case of Eq.~\eqref{eq:disc_quad_int},
\begin{widetext}
\begin{align}
    S^{(2)_{\text{int}}}={}&\frac{1}{2}\int_{t\omega xp}(\partial_{\omega}f_{0})\dot{\phi}[(\partial_{\omega}\Sigma^{R}_{0})\dot{\phi}-\Delta_{\phi}^{(1)}\Sigma^{R}]\notag\\
    ={}&\frac{1}{2}\int_{t\omega xp}\int_{\omega'\theta'}(\partial_{\omega}f_{0})\dot{\phi}D_{1}(\omega'-\omega,\theta'-\theta)(\partial_{\omega'}f_{0}')\left[\dot{\phi}-\dot{\phi}'\right]\notag\\
    ={}&\frac{1}{2}\int_{\Omega q\omega\theta}\int_{\omega'\theta'}[(\partial_{\omega}f_{0})\Omega][(\partial_{\omega'}f_{0}')\Omega]D_{1}(\omega'-\omega,\theta'-\theta)\hat{\phi}(\Omega,\mathbf{q};\omega,\theta)[\hat{\phi}(\Omega,\mathbf{q};\omega,\theta)-\hat{\phi}(\Omega,\mathbf{q};\omega',\theta')]\notag\\
    \approx{}&\frac{1}{2}\int_{\Omega q\omega\theta}\int_{\omega'\theta'}[f_{0}(\omega+\Omega)-f_{0}(\omega)][f_{0}(\omega'+\Omega)-f_{0}(\omega')]D_{1}(\omega'-\omega,\theta'-\theta)\hat{\phi}(\Omega,\mathbf{q};\omega,\theta)\notag\\
    &\qquad\qquad\qquad\times[\hat{\phi}(-\Omega,-\mathbf{q};\omega,\theta)-\hat{\phi}(-\Omega,-\mathbf{q};\omega',\theta')]\notag\\
    ={}&\frac{1}{2}\int_{\Omega q\omega\theta}\int_{\omega'\theta'}[f_{0}(\omega+\Omega)-f_{0}(\omega)][f_{0}(\omega'+\Omega)-f_{0}(\omega')]D_{1}(\omega'-\omega,\theta'-\theta)\hat{\phi}(\Omega,\mathbf{q};\omega,\theta)\hat{\phi}(-\Omega,-\mathbf{q};\omega,\theta)\notag\\
    &-\frac{1}{2}\int_{\Omega q\omega\theta}\int_{\omega'\theta'}[f_{0}(\omega+\Omega)-f_{0}(\omega)][f_{0}(\omega'+\Omega)-f_{0}(\omega')]D_{1}(\omega'-\omega,\theta'-\theta)\hat{\phi}(\Omega,\mathbf{q};\omega,\theta)\hat{\phi}(-\Omega,-\mathbf{q};\omega',\theta')\notag\\
    ={}&\frac{1}{2}\int_{\Omega q\omega\theta}\int_{\omega'\theta'}[f_{0}(\omega+\Omega)-f_{0}(\omega)][\Sigma^{R}_{0}(\omega+\Omega)-\Sigma^{R}_{0}(\omega)]\hat{\phi}(\Omega,\mathbf{q};\omega,\theta)\hat{\phi}(-\Omega,-\mathbf{q};\omega,\theta)\notag\\
    &-\frac{1}{2}\int_{\Omega q\omega\theta}\int_{\omega'\theta'}[f_{0}(\omega+\Omega)-f_{0}(\omega)][f_{0}(\omega'+\Omega)-f_{0}(\omega')]D_{1}(\omega'-\omega,\theta'-\theta)\hat{\phi}(\Omega,\mathbf{q};\omega,\theta)\hat{\phi}(-\Omega,-\mathbf{q};\omega',\theta'),
\end{align}
where we used $(\partial_{\omega}f_{0})\Omega\approx f_{0}(\omega+\Omega)-f_{0}(\omega)$.

\section{Discussion on the order of integration for the density correlation function}\label{app:int_order} 

In this section, we will examine subtleties regarding the order of integration in the evaluation of the density correlation function. A direct evaluation of the particle-hole diagram where the integration over frequency is done first yields 
\begin{align}
\Pi_{0}(i\nu,\mathbf{q})={}&-\int\frac{d\nu'}{2\pi}\frac{d^{2}p}{(2\pi)^{2}} \frac{1}{i\nu'+i\nu-\xi_{\mathbf{p}+\mathbf{q}}}\frac{1}{i\nu'-\xi_{\mathbf{p}}}\notag\\
\approx{}&N(\mu)\int\frac{d\nu'}{2\pi}\frac{d\theta d\xi_{p}}{(2\pi)^{2}}\frac{1}{i\nu-v_{F}q\cos\theta}\left[\frac{1}{i\nu'+i\nu-\xi_{p}-v_{F}q\cos\theta}-\frac{1}{i\nu'-\xi_{p}}\right]\notag\\
={}&N(\mu)\int\frac{d\theta d\xi_{p}}{(2\pi)^{2}}\frac{f_{0}(\xi_{p}+v_{F}q\cos\theta)-f_{0}(\xi_{p})}{i\nu-v_{F}q\cos\theta}\notag\\
={}&-\frac{N(\mu)}{2\pi}\int\frac{d\theta}{2\pi}\frac{v_{F}q\cos\theta}{i\nu-v_{F}q\cos\theta},
\end{align}
and after the Wick rotation, it can be written as
\begin{align}
\Pi_{0}(\Omega,\mathbf{q})={}&-\frac{N(\mu)}{2\pi}\int\frac{d\theta}{2\pi}\frac{v_{F}q\cos\theta}{\Omega-v_{F}q\cos\theta}
=\frac{N(\mu)}{2\pi}\int\frac{d\theta}{2\pi}\left[1-\frac{\Omega}{\Omega-v_{F}q\cos\theta}\right]\notag\\
={}&\frac{N(\mu)}{2\pi}\left[1-\int\frac{d\theta}{2\pi}\frac{\Omega}{\Omega-v_{F}q\cos\theta}\right]=N(\mu)\left[\frac{1}{2\pi}+\int\frac{d\omega d\theta}{(2\pi)^{2}}\frac{f_{0}(\omega+\Omega)-f_{0}(\omega)}{\Omega-v_{F}q\cos\theta}\right],
\end{align}
 which has a difference in its real part from Eq.~\eqref{eq:non-int_den_correl} in the main text. However, we can cure these apparent inconsistencies by careful consideration of the integral and introduction of a cutoff $\Lambda$ for the integration with respect to the energy. Then, doing the integral over energy first, we get
\begin{align}
\Pi_{0}(i\nu,\mathbf{q})={}&N(\mu)\int\frac{d\nu'}{2\pi}\frac{d\xi_{\mathbf{p}}d\theta}{(2\pi)^{2}}\frac{1}{i\nu-v_{F}\cos\theta}\left[\frac{1}{i\nu'+i\nu-\xi_{\mathbf{p}}-v_{F}q\cos\theta}-\frac{1}{i\nu'-\xi_{\mathbf{p}}}\right]\notag\\
={}&-N(\mu)\int\frac{d\nu'}{2\pi}\frac{d\theta}{(2\pi)^{2}}\frac{1}{i\nu-v_{F}\cos\theta}\int_{-\Lambda}^{\Lambda}d\xi_{\mathbf{p}}\left[\frac{1}{\xi_{\mathbf{p}}+v_{F}q\cos\theta-i\nu'-i\nu}-\frac{1}{\xi_{\mathbf{p}}-i\nu'}\right]\notag\\
={}&-N(\mu)\int\frac{d\nu'}{2\pi}\frac{d\theta}{(2\pi)^{2}}\frac{1}{i\nu-v_{F}\cos\theta}\left[\ln\left(\frac{\Lambda+v_{F}q\cos\theta-i\nu'-i\nu}{-\Lambda+v_{F}q\cos\theta-i\nu'-i\nu}\right)-\ln\left(\frac{\Lambda-i\nu'}{-\Lambda-i\nu'}\right)\right]\notag\\
\approx{}&-N(\mu)\int\frac{d\nu'}{2\pi}\frac{d\theta}{(2\pi)^{2}}\frac{1}{i\nu-v_{F}\cos\theta}\left[\ln\left(\frac{\Lambda-i\nu'-i\nu}{-\Lambda-i\nu'-i\nu}\right)+\frac{2\Lambda(v_{F}q\cos\theta)}{\Lambda^{2}+(\omega_{n}+\nu)^{2}}-\ln\left(\frac{\Lambda-i\nu'}{-\Lambda-i\nu'}\right)\right]\notag\\
\approx{}&-N(\mu)\int\frac{d\theta}{(2\pi)^{2}}\frac{v_{F}q\cos\theta}{i\nu-v_{F}\cos\theta}.
\end{align}
After the Wick rotation, it can be written as
\begin{align}
\Pi_{0}(\Omega,\mathbf{q})={}&-N(\mu)\int\frac{d\theta}{(2\pi)^{2}}\frac{v_{F}q\cos\theta}{\Omega-v_{F}\cos\theta}\notag\\
={}&\frac{N(\mu)}{2\pi}\left(1-\frac{\Omega}{\sqrt{\Omega^{2}-v_{F}^{2}q^{2}}}\right).
\end{align}
So we restore the well-known expression for the particle-hole bubble in two dimensions~\cite{Chubukov2003}.



\section{Comments on non-linear corrections in the general case}\label{app:general_higher}


Here, we will discuss the higher-order expansion of the action in terms of $\phi$, where we consider the full functional expansion of the self-energy. The general effective action Eq.~\eqref{eq:action} can be expanded as
\begin{align}
S={}&\int_{t\omega}\braket{f,\omega-\xi_{\mathbf{p}}-\sum_{n=1}\tfrac{1}{n+1}\Sigma^{R(n)}[f]}\notag\\
={}&\int_{t\omega}\braket{f_{0}+\Delta_{\phi}^{(1)}f+\Delta_{\phi}^{(2)}f+\cdots,\omega-\xi_{\mathbf{p}}-\sum_{n=1}\tfrac{1}{n+1}(\Sigma^{R(n)}[f_{0}]+\Delta_{\phi}^{(1)}\Sigma^{R(n)}+\Delta_{\phi}^{(2)}\Sigma^{R(n)}+\cdots)}\notag\\
={}&\int_{t\omega}\braket{f_{0},\omega-\xi_{\mathbf{p}}-\sum_{n=1}\tfrac{1}{n+1}\Sigma^{R(n)}[f_{0}]}\notag\\
&+\int_{t\omega}\braket{\Delta_{\phi}^{(1)}f,\omega-\xi_{\mathbf{p}}-\Sigma^{R}[f_{0}]}\notag\\
&+\int_{t\omega}\braket{\Delta_{\phi}^{(2)}f,\omega-\xi_{\mathbf{p}}-\Sigma^{R}[f_{0}]}-\sum_{n=1}\frac{n}{2}\int_{t\omega}\braket{\Delta_{\phi}^{(1)}f,\bar{\Delta}_{\phi}^{(1)}\Sigma^{R(n)}}\notag\\
&+\cdots.
\label{eq:general_effective_expand}
\end{align}
From Eq.~\eqref{eq:general_effective_expand}, we can obtain the general form of $m$-th order action in $\phi$ given by
\begin{align}
    S^{(m)}={}&\int_{t\omega}\braket{\Delta_{\phi}^{(m)}f,\omega-\xi_{\mathbf{p}}-\Sigma^{R}[f_{0}]}-\sum_{n=1}\sum_{\{m_{i}\}}C_{n}(\{m_{i}\})\int_{t\omega}\braket{\Delta_{\phi}^{(m_{1})}f,\bar{\Delta}_{\phi}^{(m_{2},\cdots,m_{j})}\Sigma^{R(n)}},\label{eq:general_expansion_form}
\end{align}
where $\{m_{i}\}=(m_{1},m_{2}\cdots,m_{j})$ is a partition for $m$. So we perform the summation on all the possible partitions of $m$, except $\{m_{i}\}=m$ because it is already included in the first term in Eq.~\eqref{eq:general_expansion_form}. The coefficient $C_{n}(\{m_{i}\})$ is the number of possibilities for choosing $\{m_{i}\}$ combinations divided by $(n+1)$, which is defined by $C_{n}(\{m_{i}\})=\frac{n!}{(n+1-j)!\prod_{i=1}n_{i}!}$, and $n_{i}$ is the number of $i$ in the partition and $j$ is the size of partition $\{m_{i}\}$. For example, in the case of $\int_{t\omega}\braket{\Delta_{\phi}^{(3)},\bar{\Delta}_{\phi}^{(3,2,2,1)}\Sigma^{R(n)}}$, which is one of the $\phi^{11}$ order contributions, since $n_{3}=2$, $n_{2}=2$, $n_{1}=1$, we have $C_{n}(3,3,2,2,1)=\frac{n!}{2!2!1!(n-4)!}=n(n-1)(n-2)(n-3)/4$.




As another simple example, the cubic action in $\phi$ is given by
\begin{align}
    S^{(3)}={}&\int_{t\omega}\braket{\Delta_{\phi}^{(3)}f,\omega-\xi_{p}-\Sigma^{R}[f_{0}]}-\sum_{n=1}n\int_{t\omega}\braket{\Delta_{\phi}^{(2)}f,\bar{\Delta}_{\phi}^{(1)}\Sigma^{R(n)}}
-\sum_{n=1}\frac{n(n-1)}{6}\int_{t\omega}\braket{\Delta_{\phi}^{(1)}f,\bar{\Delta}_{\phi}^{(1,1)}\Sigma^{R(n)}}.\label{eq:general_cubic}
\end{align}
The last term in Eq.~\eqref{eq:general_cubic} is completely new and not found in the one-loop order approximation (Eqs.~\eqref{eq:cont_cubic_bare} and \eqref{eq:cont_cubic_int}).
The coefficients of the second and last terms are $C_{n}(2,1)=\frac{n!}{1!1!(n-1)!}=n$ and $C_{n}(1,1,1)=\frac{n!}{3!(n-2)!}=n(n-1)/6$, respectively. That is because $n_{2}=1$, $n_{1}=1$, and $j=2$ for the second term, whereas $n_{1}=3$ and $j=3$ for the last term. For the first term, $C_{n}(3)=\frac{n!}{1!n!}=1$ since $n_{3}=1$.


\end{widetext}

%\bibliographystyle{apsrev4-1}
%\bibliography{ref}
\begin{thebibliography}{57}%
\makeatletter
\providecommand \@ifxundefined [1]{%
 \@ifx{#1\undefined}
}%
\providecommand \@ifnum [1]{%
 \ifnum #1\expandafter \@firstoftwo
 \else \expandafter \@secondoftwo
 \fi
}%
\providecommand \@ifx [1]{%
 \ifx #1\expandafter \@firstoftwo
 \else \expandafter \@secondoftwo
 \fi
}%
\providecommand \natexlab [1]{#1}%
\providecommand \enquote  [1]{``#1''}%
\providecommand \bibnamefont  [1]{#1}%
\providecommand \bibfnamefont [1]{#1}%
\providecommand \citenamefont [1]{#1}%
\providecommand \href@noop [0]{\@secondoftwo}%
\providecommand \href [0]{\begingroup \@sanitize@url \@href}%
\providecommand \@href[1]{\@@startlink{#1}\@@href}%
\providecommand \@@href[1]{\endgroup#1\@@endlink}%
\providecommand \@sanitize@url [0]{\catcode `\\12\catcode `\$12\catcode
  `\&12\catcode `\#12\catcode `\^12\catcode `\_12\catcode `\%12\relax}%
\providecommand \@@startlink[1]{}%
\providecommand \@@endlink[0]{}%
\providecommand \url  [0]{\begingroup\@sanitize@url \@url }%
\providecommand \@url [1]{\endgroup\@href {#1}{\urlprefix }}%
\providecommand \urlprefix  [0]{URL }%
\providecommand \Eprint [0]{\href }%
\providecommand \doibase [0]{http://dx.doi.org/}%
\providecommand \selectlanguage [0]{\@gobble}%
\providecommand \bibinfo  [0]{\@secondoftwo}%
\providecommand \bibfield  [0]{\@secondoftwo}%
\providecommand \translation [1]{[#1]}%
\providecommand \BibitemOpen [0]{}%
\providecommand \bibitemStop [0]{}%
\providecommand \bibitemNoStop [0]{.\EOS\space}%
\providecommand \EOS [0]{\spacefactor3000\relax}%
\providecommand \BibitemShut  [1]{\csname bibitem#1\endcsname}%
\let\auto@bib@innerbib\@empty
%</preamble>
\bibitem [{\citenamefont {Sommerfeld}(1928)}]{sommerfeld1928elektronentheorie}%
  \BibitemOpen
  \bibfield  {author} {\bibinfo {author} {\bibfnamefont {A.}~\bibnamefont
  {Sommerfeld}},\ }\href {\doibase 10.1007/BF01391052} {\bibfield  {journal}
  {\bibinfo  {journal} {Zeitschrift f{\"u}r Physik}\ }\textbf {\bibinfo
  {volume} {47}},\ \bibinfo {pages} {1} (\bibinfo {year} {1928})}\BibitemShut
  {NoStop}%
\bibitem [{\citenamefont {Bloch}(1929)}]{bloch1929quantenmechanik}%
  \BibitemOpen
  \bibfield  {author} {\bibinfo {author} {\bibfnamefont {F.}~\bibnamefont
  {Bloch}},\ }\href {\doibase 10.1007/BF01339455} {\bibfield  {journal}
  {\bibinfo  {journal} {Zeitschrift f{\"u}r physik}\ }\textbf {\bibinfo
  {volume} {52}},\ \bibinfo {pages} {555} (\bibinfo {year} {1929})}\BibitemShut
  {NoStop}%
\bibitem [{\citenamefont {Altshuler}\ \emph {et~al.}(1994)\citenamefont
  {Altshuler}, \citenamefont {Ioffe},\ and\ \citenamefont
  {Millis}}]{altshuler1994low}%
  \BibitemOpen
  \bibfield  {author} {\bibinfo {author} {\bibfnamefont {B.}~\bibnamefont
  {Altshuler}}, \bibinfo {author} {\bibfnamefont {L.}~\bibnamefont {Ioffe}}, \
  and\ \bibinfo {author} {\bibfnamefont {A.}~\bibnamefont {Millis}},\ }\href
  {\doibase 10.1103/PhysRevB.50.14048} {\bibfield  {journal} {\bibinfo
  {journal} {Physical Review B}\ }\textbf {\bibinfo {volume} {50}},\ \bibinfo
  {pages} {14048} (\bibinfo {year} {1994})}\BibitemShut {NoStop}%
\bibitem [{\citenamefont {Kim}\ \emph {et~al.}(1994)\citenamefont {Kim},
  \citenamefont {Furusaki}, \citenamefont {Wen},\ and\ \citenamefont
  {Lee}}]{kim1994gauge}%
  \BibitemOpen
  \bibfield  {author} {\bibinfo {author} {\bibfnamefont {Y.~B.}\ \bibnamefont
  {Kim}}, \bibinfo {author} {\bibfnamefont {A.}~\bibnamefont {Furusaki}},
  \bibinfo {author} {\bibfnamefont {X.-G.}\ \bibnamefont {Wen}}, \ and\
  \bibinfo {author} {\bibfnamefont {P.~A.}\ \bibnamefont {Lee}},\ }\href
  {\doibase 10.1103/PhysRevB.50.17917} {\bibfield  {journal} {\bibinfo
  {journal} {Physical Review B}\ }\textbf {\bibinfo {volume} {50}},\ \bibinfo
  {pages} {17917} (\bibinfo {year} {1994})}\BibitemShut {NoStop}%
\bibitem [{\citenamefont {Kim}\ \emph {et~al.}(1995)\citenamefont {Kim},
  \citenamefont {Lee},\ and\ \citenamefont {Wen}}]{Kim1995}%
  \BibitemOpen
  \bibfield  {author} {\bibinfo {author} {\bibfnamefont {Y.~B.}\ \bibnamefont
  {Kim}}, \bibinfo {author} {\bibfnamefont {P.~A.}\ \bibnamefont {Lee}}, \ and\
  \bibinfo {author} {\bibfnamefont {X.-G.}\ \bibnamefont {Wen}},\ }\href
  {\doibase 10.1103/PhysRevB.52.17275} {\bibfield  {journal} {\bibinfo
  {journal} {Phys. Rev. B}\ }\textbf {\bibinfo {volume} {52}},\ \bibinfo
  {pages} {17275} (\bibinfo {year} {1995})}\BibitemShut {NoStop}%
\bibitem [{\citenamefont {Polchinski}(1994)}]{polchinski1994low}%
  \BibitemOpen
  \bibfield  {author} {\bibinfo {author} {\bibfnamefont {J.}~\bibnamefont
  {Polchinski}},\ }\href {\doibase 10.1016/0550-3213(94)90449-9} {\bibfield
  {journal} {\bibinfo  {journal} {Nuclear Physics B}\ }\textbf {\bibinfo
  {volume} {422}},\ \bibinfo {pages} {617} (\bibinfo {year}
  {1994})}\BibitemShut {NoStop}%
\bibitem [{\citenamefont {Lee}\ and\ \citenamefont
  {Nagaosa}(1992)}]{lee1992gauge}%
  \BibitemOpen
  \bibfield  {author} {\bibinfo {author} {\bibfnamefont {P.~A.}\ \bibnamefont
  {Lee}}\ and\ \bibinfo {author} {\bibfnamefont {N.}~\bibnamefont {Nagaosa}},\
  }\href {\doibase 10.1103/PhysRevB.46.5621} {\bibfield  {journal} {\bibinfo
  {journal} {Physical Review B}\ }\textbf {\bibinfo {volume} {46}},\ \bibinfo
  {pages} {5621} (\bibinfo {year} {1992})}\BibitemShut {NoStop}%
\bibitem [{\citenamefont {Oganesyan}\ \emph {et~al.}(2001)\citenamefont
  {Oganesyan}, \citenamefont {Kivelson},\ and\ \citenamefont
  {Fradkin}}]{oganesyan2001quantum}%
  \BibitemOpen
  \bibfield  {author} {\bibinfo {author} {\bibfnamefont {V.}~\bibnamefont
  {Oganesyan}}, \bibinfo {author} {\bibfnamefont {S.~A.}\ \bibnamefont
  {Kivelson}}, \ and\ \bibinfo {author} {\bibfnamefont {E.}~\bibnamefont
  {Fradkin}},\ }\href {\doibase 10.1103/PhysRevB.64.195109} {\bibfield
  {journal} {\bibinfo  {journal} {Physical Review B}\ }\textbf {\bibinfo
  {volume} {64}},\ \bibinfo {pages} {195109} (\bibinfo {year}
  {2001})}\BibitemShut {NoStop}%
\bibitem [{\citenamefont {Lee}(2008)}]{lee2008stability}%
  \BibitemOpen
  \bibfield  {author} {\bibinfo {author} {\bibfnamefont {S.-S.}\ \bibnamefont
  {Lee}},\ }\href {\doibase 10.1103/PhysRevB.78.085129} {\bibfield  {journal}
  {\bibinfo  {journal} {Physical Review B}\ }\textbf {\bibinfo {volume} {78}},\
  \bibinfo {pages} {085129} (\bibinfo {year} {2008})}\BibitemShut {NoStop}%
\bibitem [{\citenamefont {Fitzpatrick}\ \emph {et~al.}(2014)\citenamefont
  {Fitzpatrick}, \citenamefont {Kachru}, \citenamefont {Kaplan},\ and\
  \citenamefont {Raghu}}]{fitzpatrick2014non}%
  \BibitemOpen
  \bibfield  {author} {\bibinfo {author} {\bibfnamefont {A.~L.}\ \bibnamefont
  {Fitzpatrick}}, \bibinfo {author} {\bibfnamefont {S.}~\bibnamefont {Kachru}},
  \bibinfo {author} {\bibfnamefont {J.}~\bibnamefont {Kaplan}}, \ and\ \bibinfo
  {author} {\bibfnamefont {S.}~\bibnamefont {Raghu}},\ }\href {\doibase 10.1103/PhysRevB.89.165114} {\bibfield  {journal} {\bibinfo  {journal}
  {Physical Review B}\ }\textbf {\bibinfo {volume} {89}},\ \bibinfo {pages}
  {165114} (\bibinfo {year} {2014})}\BibitemShut {NoStop}%
\bibitem [{\citenamefont {Abanov}\ and\ \citenamefont
  {Chubukov}(2000)}]{abanov2000spin}%
  \BibitemOpen
  \bibfield  {author} {\bibinfo {author} {\bibfnamefont {A.}~\bibnamefont
  {Abanov}}\ and\ \bibinfo {author} {\bibfnamefont {A.~V.}\ \bibnamefont
  {Chubukov}},\ }\href {\doibase 10.1103/PhysRevLett.84.5608} {\bibfield
  {journal} {\bibinfo  {journal} {Physical review letters}\ }\textbf {\bibinfo
  {volume} {84}},\ \bibinfo {pages} {5608} (\bibinfo {year}
  {2000})}\BibitemShut {NoStop}%
\bibitem [{\citenamefont {Moon}\ \emph {et~al.}(2013)\citenamefont {Moon},
  \citenamefont {Xu}, \citenamefont {Kim},\ and\ \citenamefont
  {Balents}}]{moon2013non}%
  \BibitemOpen
  \bibfield  {author} {\bibinfo {author} {\bibfnamefont {E.-G.}\ \bibnamefont
  {Moon}}, \bibinfo {author} {\bibfnamefont {C.}~\bibnamefont {Xu}}, \bibinfo
  {author} {\bibfnamefont {Y.~B.}\ \bibnamefont {Kim}}, \ and\ \bibinfo
  {author} {\bibfnamefont {L.}~\bibnamefont {Balents}},\ }\href {\doibase 10.1103/PhysRevLett.111.206401} {\bibfield  {journal} {\bibinfo  {journal}
  {Physical review letters}\ }\textbf {\bibinfo {volume} {111}},\ \bibinfo
  {pages} {206401} (\bibinfo {year} {2013})}\BibitemShut {NoStop}%
\bibitem [{\citenamefont {Han}\ \emph {et~al.}(2019)\citenamefont {Han},
  \citenamefont {Lee}, \citenamefont {Moon},\ and\ \citenamefont
  {Min}}]{han2019emergent}%
  \BibitemOpen
  \bibfield  {author} {\bibinfo {author} {\bibfnamefont {S.E.}~\bibnamefont
  {Han}}, \bibinfo {author} {\bibfnamefont {C.}~\bibnamefont {Lee}}, \bibinfo
  {author} {\bibfnamefont {E.-G.}\ \bibnamefont {Moon}}, \ and\ \bibinfo
  {author} {\bibfnamefont {H.}~\bibnamefont {Min}},\ }\href {\doibase 10.1103/PhysRevLett.122.187601} {\bibfield  {journal} {\bibinfo  {journal}
  {Physical review letters}\ }\textbf {\bibinfo {volume} {122}},\ \bibinfo
  {pages} {187601} (\bibinfo {year} {2019})}\BibitemShut {NoStop}%
\bibitem [{\citenamefont {Nayak}\ and\ \citenamefont
  {Wilczek}(1994{\natexlab{a}})}]{nayak1994non}%
  \BibitemOpen
  \bibfield  {author} {\bibinfo {author} {\bibfnamefont {C.}~\bibnamefont
  {Nayak}}\ and\ \bibinfo {author} {\bibfnamefont {F.}~\bibnamefont
  {Wilczek}},\ }\href {\doibase 10.1016/0550-3213(94)90477-4} {\bibfield
  {journal} {\bibinfo  {journal} {Nuclear Physics B}\ }\textbf {\bibinfo
  {volume} {417}},\ \bibinfo {pages} {359} (\bibinfo {year}
  {1994}{\natexlab{a}})}\BibitemShut {NoStop}%
\bibitem [{\citenamefont {Nayak}\ and\ \citenamefont
  {Wilczek}(1994{\natexlab{b}})}]{nayak1994renormalization}%
  \BibitemOpen
  \bibfield  {author} {\bibinfo {author} {\bibfnamefont {C.}~\bibnamefont
  {Nayak}}\ and\ \bibinfo {author} {\bibfnamefont {F.}~\bibnamefont
  {Wilczek}},\ }\href {\doibase 10.1016/0550-3213(94)90158-9} {\bibfield
  {journal} {\bibinfo  {journal} {Nuclear Physics B}\ }\textbf {\bibinfo
  {volume} {430}},\ \bibinfo {pages} {534} (\bibinfo {year}
  {1994}{\natexlab{b}})}\BibitemShut {NoStop}%
\bibitem [{\citenamefont {Lee}(2009)}]{lee2009low}%
  \BibitemOpen
  \bibfield  {author} {\bibinfo {author} {\bibfnamefont {S.-S.}\ \bibnamefont
  {Lee}},\ }\href {\doibase 10.1103/PhysRevB.80.165102} {\bibfield  {journal}
  {\bibinfo  {journal} {Physical Review B}\ }\textbf {\bibinfo {volume} {80}},\
  \bibinfo {pages} {165102} (\bibinfo {year} {2009})}\BibitemShut {NoStop}%
\bibitem [{\citenamefont {Metlitski}\ and\ \citenamefont
  {Sachdev}(2010{\natexlab{a}})}]{metlitski2010quantum1}%
  \BibitemOpen
  \bibfield  {author} {\bibinfo {author} {\bibfnamefont {M.~A.}\ \bibnamefont
  {Metlitski}}\ and\ \bibinfo {author} {\bibfnamefont {S.}~\bibnamefont
  {Sachdev}},\ }\href {\doibase 10.1103/PhysRevB.82.075127} {\bibfield
  {journal} {\bibinfo  {journal} {Physical Review B}\ }\textbf {\bibinfo
  {volume} {82}},\ \bibinfo {pages} {075127} (\bibinfo {year}
  {2010}{\natexlab{a}})}\BibitemShut {NoStop}%
\bibitem [{\citenamefont {Metlitski}\ and\ \citenamefont
  {Sachdev}(2010{\natexlab{b}})}]{metlitski2010quantum2}%
  \BibitemOpen
  \bibfield  {author} {\bibinfo {author} {\bibfnamefont {M.~A.}\ \bibnamefont
  {Metlitski}}\ and\ \bibinfo {author} {\bibfnamefont {S.}~\bibnamefont
  {Sachdev}},\ }\href {\doibase 10.1103/PhysRevB.82.075128} {\bibfield
  {journal} {\bibinfo  {journal} {Physical Review B}\ }\textbf {\bibinfo
  {volume} {82}},\ \bibinfo {pages} {075128} (\bibinfo {year}
  {2010}{\natexlab{b}})}\BibitemShut {NoStop}%
\bibitem [{\citenamefont {Mross}\ \emph {et~al.}(2010)\citenamefont {Mross},
  \citenamefont {McGreevy}, \citenamefont {Liu},\ and\ \citenamefont
  {Senthil}}]{mross2010controlled}%
  \BibitemOpen
  \bibfield  {author} {\bibinfo {author} {\bibfnamefont {D.~F.}\ \bibnamefont
  {Mross}}, \bibinfo {author} {\bibfnamefont {J.}~\bibnamefont {McGreevy}},
  \bibinfo {author} {\bibfnamefont {H.}~\bibnamefont {Liu}}, \ and\ \bibinfo
  {author} {\bibfnamefont {T.}~\bibnamefont {Senthil}},\ }\href {\doibase 10.1103/PhysRevB.82.045121} {\bibfield  {journal} {\bibinfo  {journal}
  {Physical Review B}\ }\textbf {\bibinfo {volume} {82}},\ \bibinfo {pages}
  {045121} (\bibinfo {year} {2010})}\BibitemShut {NoStop}%
\bibitem [{\citenamefont {Dalidovich}\ and\ \citenamefont
  {Lee}(2013)}]{dalidovich2013perturbative}%
  \BibitemOpen
  \bibfield  {author} {\bibinfo {author} {\bibfnamefont {D.}~\bibnamefont
  {Dalidovich}}\ and\ \bibinfo {author} {\bibfnamefont {S.-S.}\ \bibnamefont
  {Lee}},\ }\href {\doibase 10.1103/PhysRevB.88.245106} {\bibfield  {journal}
  {\bibinfo  {journal} {Physical Review B}\ }\textbf {\bibinfo {volume} {88}},\
  \bibinfo {pages} {245106} (\bibinfo {year} {2013})}\BibitemShut {NoStop}%
\bibitem [{\citenamefont {Lee}(2018)}]{lee2018recent}%
  \BibitemOpen
  \bibfield  {author} {\bibinfo {author} {\bibfnamefont {S.-S.}\ \bibnamefont
  {Lee}},\ }\href {\doibase 10.1146/annurev-conmatphys-031016-025531}
  {\bibfield  {journal} {\bibinfo  {journal} {Annual Review of Condensed Matter
  Physics}\ }\textbf {\bibinfo {volume} {9}},\ \bibinfo {pages} {227} (\bibinfo
  {year} {2018})}\BibitemShut {NoStop}%
\bibitem [{\citenamefont {Han}\ and\ \citenamefont {Kim}(2022)}]{han2022non}%
  \BibitemOpen
  \bibfield  {author} {\bibinfo {author} {\bibfnamefont {S.E.}~\bibnamefont
  {Han}}\ and\ \bibinfo {author} {\bibfnamefont {Y.~B.}\ \bibnamefont {Kim}},\
  }\href {\doibase 10.1103/PhysRevB.106.L081106} {\bibfield  {journal}
  {\bibinfo  {journal} {Physical Review B}\ }\textbf {\bibinfo {volume}
  {106}},\ \bibinfo {pages} {L081106} (\bibinfo {year} {2022})}\BibitemShut
  {NoStop}%
\bibitem [{\citenamefont {Coleman}\ \emph {et~al.}(2001)\citenamefont
  {Coleman}, \citenamefont {P{\'e}pin}, \citenamefont {Si},\ and\ \citenamefont
  {Ramazashvili}}]{coleman2001fermi}%
  \BibitemOpen
  \bibfield  {author} {\bibinfo {author} {\bibfnamefont {P.}~\bibnamefont
  {Coleman}}, \bibinfo {author} {\bibfnamefont {C.}~\bibnamefont {P{\'e}pin}},
  \bibinfo {author} {\bibfnamefont {Q.}~\bibnamefont {Si}}, \ and\ \bibinfo
  {author} {\bibfnamefont {R.}~\bibnamefont {Ramazashvili}},\ }\href {\doibase 10.1088/0953-8984/13/35/202} {\bibfield  {journal} {\bibinfo  {journal}
  {Journal of Physics: Condensed Matter}\ }\textbf {\bibinfo {volume} {13}},\
  \bibinfo {pages} {R723} (\bibinfo {year} {2001})}\BibitemShut {NoStop}%
\bibitem [{\citenamefont {Han}\ \emph {et~al.}(2022)\citenamefont {Han},
  \citenamefont {Schultz},\ and\ \citenamefont {Kim}}]{han2022non2}%
  \BibitemOpen
  \bibfield  {author} {\bibinfo {author} {\bibfnamefont {S.E.}~\bibnamefont
  {Han}}, \bibinfo {author} {\bibfnamefont {D.~J.}\ \bibnamefont {Schultz}}, \
  and\ \bibinfo {author} {\bibfnamefont {Y.~B.}\ \bibnamefont {Kim}},\ }\href
  {\doibase 10.1103/PhysRevB.106.155155} {\bibfield  {journal} {\bibinfo
  {journal} {Physical Review B}\ }\textbf {\bibinfo {volume} {106}},\ \bibinfo
  {pages} {155155} (\bibinfo {year} {2022})}\BibitemShut {NoStop}%
\bibitem [{\citenamefont {Schultz}\ \emph {et~al.}(2023)\citenamefont
  {Schultz}, \citenamefont {Han},\ and\ \citenamefont
  {Kim}}]{han2022microscopic}%
  \BibitemOpen
  \bibfield  {author} {\bibinfo {author} {\bibfnamefont {D.~J.}\ \bibnamefont
  {Schultz}}, \bibinfo {author} {\bibfnamefont {S.E.}~\bibnamefont {Han}}, \ and\
  \bibinfo {author} {\bibfnamefont {Y.~B.}\ \bibnamefont {Kim}},\ }\href
  {\doibase 10.1103/PhysRevB.108.L060401} {\bibfield  {journal} {\bibinfo
  {journal} {Phys. Rev. B}\ }\textbf {\bibinfo {volume} {108}},\ \bibinfo
  {pages} {L060401} (\bibinfo {year} {2023})}\BibitemShut {NoStop}%
\bibitem [{\citenamefont {Senthil}\ \emph {et~al.}(2004)\citenamefont
  {Senthil}, \citenamefont {Vojta},\ and\ \citenamefont
  {Sachdev}}]{senthil2004weak}%
  \BibitemOpen
  \bibfield  {author} {\bibinfo {author} {\bibfnamefont {T.}~\bibnamefont
  {Senthil}}, \bibinfo {author} {\bibfnamefont {M.}~\bibnamefont {Vojta}}, \
  and\ \bibinfo {author} {\bibfnamefont {S.}~\bibnamefont {Sachdev}},\ }\href
  {\doibase 10.1103/PhysRevB.69.035111} {\bibfield  {journal} {\bibinfo
  {journal} {Physical Review B}\ }\textbf {\bibinfo {volume} {69}},\ \bibinfo
  {pages} {035111} (\bibinfo {year} {2004})}\BibitemShut {NoStop}%
\bibitem [{\citenamefont {Senthil}(2008{\natexlab{a}})}]{senthil2008critical}%
  \BibitemOpen
  \bibfield  {author} {\bibinfo {author} {\bibfnamefont {T.}~\bibnamefont
  {Senthil}},\ }\href {\doibase 10.1103/PhysRevB.78.035103} {\bibfield
  {journal} {\bibinfo  {journal} {Physical Review B}\ }\textbf {\bibinfo
  {volume} {78}},\ \bibinfo {pages} {035103} (\bibinfo {year}
  {2008}{\natexlab{a}})}\BibitemShut {NoStop}%
\bibitem [{\citenamefont {Senthil}(2008{\natexlab{b}})}]{senthil2008theory}%
  \BibitemOpen
  \bibfield  {author} {\bibinfo {author} {\bibfnamefont {T.}~\bibnamefont
  {Senthil}},\ }\href {\doibase 10.1103/PhysRevB.78.045109} {\bibfield
  {journal} {\bibinfo  {journal} {Physical Review B}\ }\textbf {\bibinfo
  {volume} {78}},\ \bibinfo {pages} {045109} (\bibinfo {year}
  {2008}{\natexlab{b}})}\BibitemShut {NoStop}%
\bibitem [{\citenamefont {Sachdev}\ and\ \citenamefont
  {Ye}(1993)}]{sachdev1993gapless}%
  \BibitemOpen
  \bibfield  {author} {\bibinfo {author} {\bibfnamefont {S.}~\bibnamefont
  {Sachdev}}\ and\ \bibinfo {author} {\bibfnamefont {J.}~\bibnamefont {Ye}},\
  }\href {\doibase 10.1103/PhysRevLett.70.3339} {\bibfield  {journal} {\bibinfo
   {journal} {Physical review letters}\ }\textbf {\bibinfo {volume} {70}},\
  \bibinfo {pages} {3339} (\bibinfo {year} {1993})}\BibitemShut {NoStop}%
\bibitem [{\citenamefont {Sachdev}(2010)}]{sachdev2010holographic}%
  \BibitemOpen
  \bibfield  {author} {\bibinfo {author} {\bibfnamefont {S.}~\bibnamefont
  {Sachdev}},\ }\href {\doibase 10.1103/PhysRevLett.105.151602} {\bibfield
  {journal} {\bibinfo  {journal} {Physical review letters}\ }\textbf {\bibinfo
  {volume} {105}},\ \bibinfo {pages} {151602} (\bibinfo {year}
  {2010})}\BibitemShut {NoStop}%
\bibitem [{\citenamefont {Maldacena}\ and\ \citenamefont
  {Stanford}(2016)}]{maldacena2016remarks}%
  \BibitemOpen
  \bibfield  {author} {\bibinfo {author} {\bibfnamefont {J.}~\bibnamefont
  {Maldacena}}\ and\ \bibinfo {author} {\bibfnamefont {D.}~\bibnamefont
  {Stanford}},\ }\href {\doibase 10.1103/PhysRevD.94.106002} {\bibfield
  {journal} {\bibinfo  {journal} {Physical Review D}\ }\textbf {\bibinfo
  {volume} {94}},\ \bibinfo {pages} {106002} (\bibinfo {year}
  {2016})}\BibitemShut {NoStop}%
\bibitem [{\citenamefont {Chowdhury}\ \emph {et~al.}(2022)\citenamefont
  {Chowdhury}, \citenamefont {Georges}, \citenamefont {Parcollet},\ and\
  \citenamefont {Sachdev}}]{chowdhury2022sachdev}%
  \BibitemOpen
  \bibfield  {author} {\bibinfo {author} {\bibfnamefont {D.}~\bibnamefont
  {Chowdhury}}, \bibinfo {author} {\bibfnamefont {A.}~\bibnamefont {Georges}},
  \bibinfo {author} {\bibfnamefont {O.}~\bibnamefont {Parcollet}}, \ and\
  \bibinfo {author} {\bibfnamefont {S.}~\bibnamefont {Sachdev}},\ }\href
  {\doibase 10.1103/RevModPhys.94.035004} {\bibfield  {journal} {\bibinfo
  {journal} {Reviews of Modern Physics}\ }\textbf {\bibinfo {volume} {94}},\
  \bibinfo {pages} {035004} (\bibinfo {year} {2022})}\BibitemShut {NoStop}%
\bibitem [{\citenamefont {Shi}\ \emph {et~al.}(2022)\citenamefont {Shi},
  \citenamefont {Goldman}, \citenamefont {Else},\ and\ \citenamefont
  {Senthil}}]{shi2022gifts}%
  \BibitemOpen
  \bibfield  {author} {\bibinfo {author} {\bibfnamefont {Z.~D.}\ \bibnamefont
  {Shi}}, \bibinfo {author} {\bibfnamefont {H.}~\bibnamefont {Goldman}},
  \bibinfo {author} {\bibfnamefont {D.~V.}\ \bibnamefont {Else}}, \ and\
  \bibinfo {author} {\bibfnamefont {T.}~\bibnamefont {Senthil}},\ }\href
  {\doibase 10.21468/SciPostPhys.13.5.102} {\bibfield  {journal} {\bibinfo
  {journal} {SciPost Physics}\ }\textbf {\bibinfo {volume} {13}},\ \bibinfo
  {pages} {102} (\bibinfo {year} {2022})}\BibitemShut {NoStop}%
\bibitem [{\citenamefont {Else}\ \emph {et~al.}(2021)\citenamefont {Else},
  \citenamefont {Thorngren},\ and\ \citenamefont {Senthil}}]{else2021non}%
  \BibitemOpen
  \bibfield  {author} {\bibinfo {author} {\bibfnamefont {D.~V.}\ \bibnamefont
  {Else}}, \bibinfo {author} {\bibfnamefont {R.}~\bibnamefont {Thorngren}}, \
  and\ \bibinfo {author} {\bibfnamefont {T.}~\bibnamefont {Senthil}},\ }\href
  {\doibase 10.1103/PhysRevX.11.021005} {\bibfield  {journal} {\bibinfo
  {journal} {Physical Review X}\ }\textbf {\bibinfo {volume} {11}},\ \bibinfo
  {pages} {021005} (\bibinfo {year} {2021})}\BibitemShut {NoStop}%
\bibitem [{\citenamefont {Else}\ and\ \citenamefont
  {Senthil}(2021)}]{else2021strange}%
  \BibitemOpen
  \bibfield  {author} {\bibinfo {author} {\bibfnamefont {D.~V.}\ \bibnamefont
  {Else}}\ and\ \bibinfo {author} {\bibfnamefont {T.}~\bibnamefont {Senthil}},\
  }\href {\doibase 10.1103/PhysRevLett.127.086601} {\bibfield  {journal}
  {\bibinfo  {journal} {Physical Review Letters}\ }\textbf {\bibinfo {volume}
  {127}},\ \bibinfo {pages} {086601} (\bibinfo {year} {2021})}\BibitemShut
  {NoStop}%
\bibitem [{\citenamefont {Lu}\ \emph {et~al.}(2023)\citenamefont {Lu},
  \citenamefont {Wang},\ and\ \citenamefont {You}}]{lu2023definition}%
  \BibitemOpen
  \bibfield  {author} {\bibinfo {author} {\bibfnamefont {D.-C.}\ \bibnamefont
  {Lu}}, \bibinfo {author} {\bibfnamefont {J.}~\bibnamefont {Wang}}, \ and\
  \bibinfo {author} {\bibfnamefont {Y.-Z.}\ \bibnamefont {You}},\ }\href
  {\doibase 10.48550/arXiv.2302.12731} {\bibfield  {journal} {\bibinfo
  {journal} {arXiv preprint arXiv:2302.12731}\ } (\bibinfo {year} {2023}),\
  10.48550/arXiv.2302.12731}\BibitemShut {NoStop}%
\bibitem [{\citenamefont {Else}(2023)}]{else2023collisionless}%
  \BibitemOpen
  \bibfield  {author} {\bibinfo {author} {\bibfnamefont {D.~V.}\ \bibnamefont
  {Else}},\ }\href {\doibase 10.1103/PhysRevB.108.045107} {\bibfield  {journal}
  {\bibinfo  {journal} {Phys. Rev. B}\ }\textbf {\bibinfo {volume} {108}},\
  \bibinfo {pages} {045107} (\bibinfo {year} {2023})}\BibitemShut {NoStop}%
\bibitem [{\citenamefont {Shi}\ \emph {et~al.}(2023)\citenamefont {Shi},
  \citenamefont {Else}, \citenamefont {Goldman},\ and\ \citenamefont
  {Senthil}}]{shi2023loop}%
  \BibitemOpen
  \bibfield  {author} {\bibinfo {author} {\bibfnamefont {Z.~D.}\ \bibnamefont
  {Shi}}, \bibinfo {author} {\bibfnamefont {D.~V.}\ \bibnamefont {Else}},
  \bibinfo {author} {\bibfnamefont {H.}~\bibnamefont {Goldman}}, \ and\
  \bibinfo {author} {\bibfnamefont {T.}~\bibnamefont {Senthil}},\ }\href
  {\doibase 10.21468/SciPostPhys.14.5.113} {\bibfield  {journal} {\bibinfo
  {journal} {SciPost Physics}\ }\textbf {\bibinfo {volume} {14}},\ \bibinfo
  {pages} {113} (\bibinfo {year} {2023})}\BibitemShut {NoStop}%
\bibitem [{\citenamefont {Delacr\'etaz}\ \emph {et~al.}(2022)\citenamefont
  {Delacr\'etaz}, \citenamefont {Du}, \citenamefont {Mehta},\ and\
  \citenamefont {Son}}]{Son2022}%
  \BibitemOpen
  \bibfield  {author} {\bibinfo {author} {\bibfnamefont {L.~V.}\ \bibnamefont
  {Delacr\'etaz}}, \bibinfo {author} {\bibfnamefont {Y.-H.}\ \bibnamefont
  {Du}}, \bibinfo {author} {\bibfnamefont {U.}~\bibnamefont {Mehta}}, \ and\
  \bibinfo {author} {\bibfnamefont {D.~T.}\ \bibnamefont {Son}},\ }\href
  {\doibase 10.1103/PhysRevResearch.4.033131} {\bibfield  {journal} {\bibinfo
  {journal} {Phys. Rev. Res.}\ }\textbf {\bibinfo {volume} {4}},\ \bibinfo
  {pages} {033131} (\bibinfo {year} {2022})}\BibitemShut {NoStop}%
\bibitem [{\citenamefont {Mandal}(2022)}]{Mandal2022}%
  \BibitemOpen
  \bibfield  {author} {\bibinfo {author} {\bibfnamefont {I.}~\bibnamefont
  {Mandal}},\ }\href {\doibase https://doi.org/10.1016/j.physleta.2022.128292}
  {\bibfield  {journal} {\bibinfo  {journal} {Physics Letters A}\ }\textbf
  {\bibinfo {volume} {447}},\ \bibinfo {pages} {128292} (\bibinfo {year}
  {2022})}\BibitemShut {NoStop}%
\bibitem [{\citenamefont {Kita}(2010)}]{kita2010introduction}%
  \BibitemOpen
  \bibfield  {author} {\bibinfo {author} {\bibfnamefont {T.}~\bibnamefont
  {Kita}},\ }\href {\doibase 10.1143/PTP.123.581} {\bibfield  {journal}
  {\bibinfo  {journal} {Progress of theoretical physics}\ }\textbf {\bibinfo
  {volume} {123}},\ \bibinfo {pages} {581} (\bibinfo {year}
  {2010})}\BibitemShut {NoStop}%
\bibitem [{\citenamefont {Rammer}\ and\ \citenamefont
  {Smith}(1986)}]{rammer1986quantum}%
  \BibitemOpen
  \bibfield  {author} {\bibinfo {author} {\bibfnamefont {J.}~\bibnamefont
  {Rammer}}\ and\ \bibinfo {author} {\bibfnamefont {H.}~\bibnamefont {Smith}},\
  }\href {\doibase 10.1103/RevModPhys.58.323} {\bibfield  {journal} {\bibinfo
  {journal} {Reviews of modern physics}\ }\textbf {\bibinfo {volume} {58}},\
  \bibinfo {pages} {323} (\bibinfo {year} {1986})}\BibitemShut {NoStop}%
\bibitem [{\citenamefont {Rammer}(2011)}]{rammer2011quantum}%
  \BibitemOpen
  \bibfield  {author} {\bibinfo {author} {\bibfnamefont {J.}~\bibnamefont
  {Rammer}},\ }\href {\doibase 10.1017/CBO9780511618956} {\bibfield  {journal}
  {\bibinfo  {journal} {Quantum Field Theory of Non-equilibrium States}\ }
  (\bibinfo {year} {2011}),\ 10.1017/CBO9780511618956}\BibitemShut {NoStop}%
\bibitem [{\citenamefont {Bruus}\ and\ \citenamefont
  {Flensberg}(2004)}]{bruus2004many}%
  \BibitemOpen
  \bibfield  {author} {\bibinfo {author} {\bibfnamefont {H.}~\bibnamefont
  {Bruus}}\ and\ \bibinfo {author} {\bibfnamefont {K.}~\bibnamefont
  {Flensberg}},\ }\href@noop {} {\emph {\bibinfo {title} {Many-body quantum
  theory in condensed matter physics: an introduction}}}\ (\bibinfo
  {publisher} {OUP Oxford},\ \bibinfo {year} {2004})\BibitemShut {NoStop}%
\bibitem [{\citenamefont {Negele}(2018)}]{negele2018quantum}%
  \BibitemOpen
  \bibfield  {author} {\bibinfo {author} {\bibfnamefont {J.~W.}\ \bibnamefont
  {Negele}},\ }\href@noop {} {\emph {\bibinfo {title} {Quantum many-particle
  systems}}}\ (\bibinfo  {publisher} {CRC Press},\ \bibinfo {year}
  {2018})\BibitemShut {NoStop}%
\bibitem [{\citenamefont {Mahan}(2000)}]{mahan2000many}%
  \BibitemOpen
  \bibfield  {author} {\bibinfo {author} {\bibfnamefont {G.~D.}\ \bibnamefont
  {Mahan}},\ }\href@noop {} {\emph {\bibinfo {title} {Many-particle physics}}}\
  (\bibinfo  {publisher} {Springer Science \& Business Media},\ \bibinfo {year}
  {2000})\BibitemShut {NoStop}%
\bibitem [{\citenamefont {Prange}\ and\ \citenamefont
  {Kadanoff}(1964)}]{Kadanoff1964}%
  \BibitemOpen
  \bibfield  {author} {\bibinfo {author} {\bibfnamefont {R.~E.}\ \bibnamefont
  {Prange}}\ and\ \bibinfo {author} {\bibfnamefont {L.~P.}\ \bibnamefont
  {Kadanoff}},\ }\href {\doibase 10.1103/PhysRev.134.A566} {\bibfield
  {journal} {\bibinfo  {journal} {Phys. Rev.}\ }\textbf {\bibinfo {volume}
  {134}},\ \bibinfo {pages} {A566} (\bibinfo {year} {1964})}\BibitemShut
  {NoStop}%
\bibitem [{\citenamefont {Hall}\ and\ \citenamefont
  {Hall}(2013)}]{hall2013lie}%
  \BibitemOpen
  \bibfield  {author} {\bibinfo {author} {\bibfnamefont {B.~C.}\ \bibnamefont
  {Hall}}\ and\ \bibinfo {author} {\bibfnamefont {B.~C.}\ \bibnamefont
  {Hall}},\ }\href {\doibase 10.1007/978-3-319-13467-3} {\emph {\bibinfo
  {title} {Lie groups, Lie algebras, and representations}}}\ (\bibinfo
  {publisher} {Springer},\ \bibinfo {year} {2013})\BibitemShut {NoStop}%
\bibitem [{\citenamefont {Miller}(1973)}]{miller1973symmetry}%
  \BibitemOpen
  \bibfield  {author} {\bibinfo {author} {\bibfnamefont {W.}~\bibnamefont
  {Miller}},\ }\href@noop {} {\emph {\bibinfo {title} {Symmetry groups and
  their applications}}}\ (\bibinfo  {publisher} {Academic Press},\ \bibinfo
  {year} {1973})\BibitemShut {NoStop}%
\bibitem [{\citenamefont {Luttinger}\ and\ \citenamefont
  {Ward}(1960)}]{luttinger1960ground}%
  \BibitemOpen
  \bibfield  {author} {\bibinfo {author} {\bibfnamefont {J.~M.}\ \bibnamefont
  {Luttinger}}\ and\ \bibinfo {author} {\bibfnamefont {J.~C.}\ \bibnamefont
  {Ward}},\ }\href {\doibase 10.1103/PhysRev.118.1417} {\bibfield  {journal}
  {\bibinfo  {journal} {Physical Review}\ }\textbf {\bibinfo {volume} {118}},\
  \bibinfo {pages} {1417} (\bibinfo {year} {1960})}\BibitemShut {NoStop}%
\bibitem [{\citenamefont {Baym}\ and\ \citenamefont
  {Kadanoff}(1961)}]{baym1961conservation}%
  \BibitemOpen
  \bibfield  {author} {\bibinfo {author} {\bibfnamefont {G.}~\bibnamefont
  {Baym}}\ and\ \bibinfo {author} {\bibfnamefont {L.~P.}\ \bibnamefont
  {Kadanoff}},\ }\href {\doibase 10.1103/PhysRev.124.287} {\bibfield  {journal}
  {\bibinfo  {journal} {Physical Review}\ }\textbf {\bibinfo {volume} {124}},\
  \bibinfo {pages} {287} (\bibinfo {year} {1961})}\BibitemShut {NoStop}%
\bibitem [{\citenamefont {Kotliar}\ \emph {et~al.}(2006)\citenamefont
  {Kotliar}, \citenamefont {Savrasov}, \citenamefont {Haule}, \citenamefont
  {Oudovenko}, \citenamefont {Parcollet},\ and\ \citenamefont
  {Marianetti}}]{kotliar2006electronic}%
  \BibitemOpen
  \bibfield  {author} {\bibinfo {author} {\bibfnamefont {G.}~\bibnamefont
  {Kotliar}}, \bibinfo {author} {\bibfnamefont {S.~Y.}\ \bibnamefont
  {Savrasov}}, \bibinfo {author} {\bibfnamefont {K.}~\bibnamefont {Haule}},
  \bibinfo {author} {\bibfnamefont {V.~S.}\ \bibnamefont {Oudovenko}}, \bibinfo
  {author} {\bibfnamefont {O.}~\bibnamefont {Parcollet}}, \ and\ \bibinfo
  {author} {\bibfnamefont {C.}~\bibnamefont {Marianetti}},\ }\href {\doibase 10.1103/RevModPhys.78.865} {\bibfield  {journal} {\bibinfo  {journal}
  {Reviews of Modern Physics}\ }\textbf {\bibinfo {volume} {78}},\ \bibinfo
  {pages} {865} (\bibinfo {year} {2006})}\BibitemShut {NoStop}%
\bibitem [{\citenamefont {Potthoff}(2004)}]{potthoff2004non}%
  \BibitemOpen
  \bibfield  {author} {\bibinfo {author} {\bibfnamefont {M.}~\bibnamefont
  {Potthoff}},\ }\href@noop {} {\bibfield  {journal} {\bibinfo  {journal}
  {arXiv preprint cond-mat/0406671}\ } (\bibinfo {year} {2004})}\BibitemShut
  {NoStop}%
\bibitem [{\citenamefont {Kozik}\ \emph {et~al.}(2015)\citenamefont {Kozik},
  \citenamefont {Ferrero},\ and\ \citenamefont
  {Georges}}]{kozik2015nonexistence}%
  \BibitemOpen
  \bibfield  {author} {\bibinfo {author} {\bibfnamefont {E.}~\bibnamefont
  {Kozik}}, \bibinfo {author} {\bibfnamefont {M.}~\bibnamefont {Ferrero}}, \
  and\ \bibinfo {author} {\bibfnamefont {A.}~\bibnamefont {Georges}},\ }\href
  {\doibase 10.1103/PhysRevLett.114.156402} {\bibfield  {journal} {\bibinfo
  {journal} {Phys. Rev. Lett.}\ }\textbf {\bibinfo {volume} {114}},\ \bibinfo
  {pages} {156402} (\bibinfo {year} {2015})}\BibitemShut {NoStop}%
\bibitem [{\citenamefont {Moyal}(1949)}]{moyal1949quantum}%
  \BibitemOpen
  \bibfield  {author} {\bibinfo {author} {\bibfnamefont {J.~E.}\ \bibnamefont
  {Moyal}},\ }in\ \href {\doibase 10.1017/S0305004100000487} {\emph {\bibinfo
  {booktitle} {Mathematical Proceedings of the Cambridge Philosophical
  Society}}},\ Vol.~\bibinfo {volume} {45}\ (\bibinfo {organization} {Cambridge
  University Press},\ \bibinfo {year} {1949})\ pp.\ \bibinfo {pages}
  {99--124}\BibitemShut {NoStop}%
\bibitem [{\citenamefont {Groenewold}(1946)}]{groenewold1946principles}%
  \BibitemOpen
  \bibfield  {author} {\bibinfo {author} {\bibfnamefont {H.~J.}\ \bibnamefont
  {Groenewold}},\ }\href {\doibase 10.1016/S0031-8914(46)80059-4} {\emph
  {\bibinfo {title} {On the principles of elementary quantum mechanics}}}\
  (\bibinfo  {publisher} {Springer},\ \bibinfo {year} {1946})\BibitemShut
  {NoStop}%
\bibitem [{\citenamefont {Chubukov}\ and\ \citenamefont
  {Maslov}(2003)}]{Chubukov2003}%
  \BibitemOpen
  \bibfield  {author} {\bibinfo {author} {\bibfnamefont {A.~V.}\ \bibnamefont
  {Chubukov}}\ and\ \bibinfo {author} {\bibfnamefont {D.~L.}\ \bibnamefont
  {Maslov}},\ }\href {\doibase 10.1103/PhysRevB.68.155113} {\bibfield
  {journal} {\bibinfo  {journal} {Phys. Rev. B}\ }\textbf {\bibinfo {volume}
  {68}},\ \bibinfo {pages} {155113} (\bibinfo {year} {2003})}\BibitemShut
  {NoStop}%
\end{thebibliography}%
\end{document}